\documentclass[11pt]{article}

\usepackage[margin=1in]{geometry}
\usepackage[T1]{fontenc}
\usepackage[utf8]{inputenc}
\usepackage{lmodern}
\usepackage{microtype}
\usepackage{amsmath,amssymb,amsthm}
\usepackage{mathtools}
\usepackage{hyperref}
\usepackage{booktabs}
\usepackage{enumitem}
\usepackage{listings}
\usepackage{xcolor}

\hypersetup{
  colorlinks=true,
  linkcolor=blue!50!black,
  citecolor=blue!50!black,
  urlcolor=blue!50!black
}

\lstdefinelanguage{Lean4}{
  morekeywords={
    import,namespace,end,open,section,noncomputable,abbrev,
    def,theorem,lemma,where,by,fun,let,in,have,show,exact,
    intro,apply,simpa,calc
  },
  sensitive=true,
  morecomment=[l]{--},
  morecomment=[s]{/-}{-/},
  morestring=[b]"
}

\lstset{
  language=Lean4,
  basicstyle=\ttfamily\small,
  keywordstyle=\bfseries\color{blue!50!black},
  commentstyle=\itshape\color{green!40!black},
  stringstyle=\color{red!50!black},
  showstringspaces=false,
  breaklines=true,
  frame=single,
  framerule=0.3pt,
  xleftmargin=0.5em,
  xrightmargin=0.5em,
  keepspaces=true,
  columns=fullflexible,
  literate=
    {¬}{{$\neg$}}1
    {∃}{{$\exists$}}1
    {→}{{$\to$}}1
    {≤}{{$\leq$}}1
}

\newtheorem{theorem}{Theorem}[section]
\newtheorem{lemma}[theorem]{Lemma}
\theoremstyle{remark}
\newtheorem{remark}[theorem]{Remark}

\title{Navier--Stokes Grassroots:\\
Ground-Up Background Theory for the Goertzel SG-Cole-Hopf Program}
\author{}
\date{\today}

\begin{document}
\maketitle

\begin{abstract}
This note records the \emph{grassroots} side of the current work around Ben
Goertzel's December 11, 2025 Navier--Stokes manuscript.  It is intentionally
not the blocker note.  The companion crux-first analysis lives separately in
\texttt{ns\_crux.tex}.

The goal here is narrower and more constructive: isolate reusable background
theory that would still matter even if the SG-Cole-Hopf route needs revision.
The first machine-checked kernel now formalizes the exact algebraic inequalities
behind the manuscript's identity-energy route: finite carré-du-champ energy,
log-gradient reduction under a positive lower bound, the \(-2\nu \log\) scaling
factor, and the Cauchy--Schwarz push-down from coefficient energy at the
identity to pointwise vorticity/frame bounds.  It now also includes an abstract
Cole--Hopf identity interface: if one can supply uniform positivity, an
identity-energy bound for \(P_t\phi\), and a frame-curl bound, the whole
algebraic chain to a uniform vorticity estimate is derived in Lean.  And it now
adds a local kernel-founded semigroup shell on top of Mettapedia's own
Markov-kernel infrastructure, so the next NS layer can stay inside the local
library rather than reaching outward prematurely.  The newest pass also adds a
formal positivity gate for this log layer: a heat field with a zero or
nonpositive value cannot instantiate the Cole--Hopf identity interface at all,
so any nonnegative-heat repair must be strengthened back to a uniform strict
positive lower bound before the existing log-energy estimates apply.  This has
now been sharpened as well: even a pointwise-positive heat field cannot enter
the interface if its values are arbitrarily close to zero, and if a scalar
derivative stays uniformly nonzero while \(\Phi\) approaches zero then the
quotient \(d\Phi/\Phi\) and the associated one-direction log energy are
unbounded.  Conversely, a uniformly bounded quotient on such a field now
formally forces the derivative numerator itself to have arbitrarily small
absolute values.  The positive
approximation shell says that if a future finite-mode or truncation theorem
shows convergence of the manuscript's radius and statistic observables, then
the induced cutoff potential and Cole--Hopf field convergence now follow
automatically, and any eventual absolute truncation bound on those quantities
passes to the limit as well.  On the identity-energy side, a new finite-mode
limit shell now says that eventual truncation positivity and eventual
truncation coefficient-energy bounds promote to honest limiting
\texttt{ColeHopfIdentityData}, so the existing vorticity bridge can be reused on
the limit object.  And the same finite-mode transfer layer now fuses with the
local kernel-semigroup shell, so the kernel-founded NS interface can also be
recovered directly on the limit data.  The newest pass packages the whole
manuscript-shaped truncation target in one Lean structure: chart-observable
convergence, cutoff/Cole--Hopf convergence, finite-mode positivity and
identity-energy transfer, and kernel-founded vorticity consequences now sit in
one explicit object.  And the newest concrete step lands the D.4 chart side in
an actual finite-mode model: a squared weighted chart radius, a bounded linear
observable, the chart-cutoff datum \(W_0 = \chi \cdot \ell\), the heat datum
\(\phi = \exp(-W_0/(2\nu))\), explicit \(W_0/\phi\) bounds, and truncation
convergence for both \(W_0\) and \(\phi\) are now machine-checked in one local
module.  The latest pass then packages the first genuinely Cole--Hopf-shaped
shared truncation instance into direct named kernel and vorticity objects, so
the cutoff-convergence and uniform vorticity consequences can now be cited from
one concrete package instead of from the whole abstraction stack.  The newest
step then replaces the static selected-coefficient \(d\Phi\) surrogate by a
nonnegative-time heat-decayed one, while preserving the same shared
truncation, positivity, and coefficient-energy consequences.  The newest pass
then tests a genuine smooth window cutoff rather than another monotone
saturation.  Lean now shows that this window still passes through the present
approximation shell and still satisfies the quantitative boundedness hypothesis
needed by the heat-decayed package.  On top of that concrete windowed package,
the route now fixes actual candidate fields, a seeded local target, a scaled
descendant bridge, and an affine seed-blend descendant bridge, together with
same-route rigidity theorems showing exactly when those descendants can still
collapse back to original self-compatibility.

This does \emph{not} prove 3D Navier--Stokes regularity.  What it does is make
the positive architecture cleaner: if the hard analytic bridges can be supplied,
the identity-energy argument now lands in a small Lean-checked interface rather
than in manuscript prose alone.
\end{abstract}

\tableofcontents

\section{Scope}

The December 2025 manuscript
\emph{``Global Regularity of 3D Navier--Stokes via SG-Cole-Hopf Program''}
proposes a route with four main layers:
\begin{enumerate}[leftmargin=1.5em]
  \item a polynomially weighted divergence-free Fourier frame on the torus;
  \item an SG diffusion and Cole--Hopf transform on the diffeomorphism group;
  \item an identity-based push-down back to Eulerian velocity fields;
  \item a first-order energy estimate at the identity strong enough to control
  vorticity and close the BKM criterion.
\end{enumerate}

The companion crux note asks where this route may fail.  This note asks a
different question:
\begin{quote}
What mathematical infrastructure do we want in Lean even if the final
regularity proof still needs serious repair?
\end{quote}

That changes the work style.  We prefer small exact interfaces over speculative
closure.  In particular:
\begin{enumerate}[leftmargin=1.5em]
  \item separate algebraic identities from heavy infinite-dimensional analysis;
  \item formalize the inequalities the manuscript repeatedly reuses;
  \item make the hard analytic inputs explicit, so future repairs can plug into
  them rather than hiding them.
\end{enumerate}

There is also a build-sanity rule now:
\begin{quote}
build from Mettapedia plus the locally pinned \texttt{v4.28.0} mathlib first,
and treat outside Lean repos as reference material only unless a direct local
need appears.
\end{quote}

\section{Existing Base}

The grassroots route is not starting from zero.  The project already has:
\begin{center}
\begin{tabular}{p{0.46\textwidth}p{0.42\textwidth}}
\toprule
\textbf{Module} & \textbf{Grassroots role} \\
\midrule
\texttt{Mathlib} finite sums / real inequalities & Cauchy--Schwarz and square-root algebra \\
\shortstack[l]{\texttt{Mettapedia/ProbabilityTheory/}\\\texttt{MarkovCategory/Kernels.lean}} & stable local kernel semantics \\
\shortstack[l]{\texttt{Mettapedia/FluidDynamics/}\\\texttt{NavierStokes/IdentityEnergy.lean}} & new NS identity-energy kernel \\
\shortstack[l]{\texttt{Mettapedia/FluidDynamics/}\\\texttt{NavierStokes/ColeHopfPositiveLowerBound.lean}} & strict positivity gate and quotient blow-up for the Cole--Hopf log transform \\
\shortstack[l]{\texttt{Mettapedia/FluidDynamics/}\\\texttt{NavierStokes/ColeHopfInterface.lean}} & abstract manuscript-shaped identity bridge \\
\shortstack[l]{\texttt{Mettapedia/FluidDynamics/}\\\texttt{NavierStokes/KernelSemigroupInterface.lean}} & local kernel-founded semigroup shell \\
\shortstack[l]{\texttt{Mettapedia/FluidDynamics/}\\\texttt{NavierStokes/CutoffContinuity.lean}} & generic cutoff continuity / instability shell \\
\shortstack[l]{\texttt{Mettapedia/FluidDynamics/}\\\texttt{NavierStokes/ColeHopfTopology.lean}} & local \(W \mapsto \phi\) continuity / instability shell \\
\shortstack[l]{\texttt{Mettapedia/FluidDynamics/}\\\texttt{NavierStokes/ApproximationInterface.lean}} & positive truncation / approximation shell \\
\shortstack[l]{\texttt{Mettapedia/FluidDynamics/}\\\texttt{NavierStokes/ManuscriptChartDatum.lean}} & concrete D.4-style chart datum on a finite-mode model \\
\shortstack[l]{\texttt{Mettapedia/FluidDynamics/}\\\texttt{NavierStokes/ManuscriptChartPackage.lean}} & concrete chart datum plugged into the combined truncation package \\
\shortstack[l]{\texttt{Mettapedia/FluidDynamics/}\\\texttt{NavierStokes/SharedApproximationPackage.lean}} & one shared `approx` for chart and identity lanes \\
\shortstack[l]{\texttt{Mettapedia/FluidDynamics/}\\\texttt{NavierStokes/GeometricSharedApproximation.lean}} & first explicit shared truncation instance with positivity and energy bounds \\
\shortstack[l]{\texttt{Mettapedia/FluidDynamics/}\\\texttt{NavierStokes/GeometricColeHopfSharedApproximation.lean}} & Cole--Hopf-shaped shared instance driven by the actual chart heat datum \\
\shortstack[l]{\texttt{Mettapedia/FluidDynamics/}\\\texttt{NavierStokes/GeometricColeHopfPackage.lean}} & direct concrete package-level kernel and vorticity consequences for that instance \\
\shortstack[l]{\texttt{Mettapedia/FluidDynamics/}\\\texttt{NavierStokes/GeometricColeHopfHeatApproximation.lean}} & nonnegative-time heat-decayed \(d\Phi\) model on the same shared truncation family \\
\shortstack[l]{\texttt{Mettapedia/FluidDynamics/}\\\texttt{NavierStokes/WindowedCutoffApproximation.lean}} & genuine smooth window still compatible with the approximation shell \\
\shortstack[l]{\texttt{Mettapedia/FluidDynamics/}\\\texttt{NavierStokes/WindowedColeHopfHeatPackage.lean}} & concrete windowed heat package and uniform vorticity tendril \\
\shortstack[l]{\texttt{Mettapedia/FluidDynamics/}\\\texttt{NavierStokes/WindowedColeHopfHeatFoundationFrontier.lean}} & fixed candidate fields and compatibility frontier on the windowed route \\
\shortstack[l]{\texttt{Mettapedia/FluidDynamics/}\\\texttt{NavierStokes/WindowedColeHopfHeatSeededModel.lean}} & stronger zero-pressure / seeded-slice local target on that route \\
\shortstack[l]{\texttt{Mettapedia/FluidDynamics/}\\\texttt{NavierStokes/WindowedColeHopfHeatScaledSeededFrontier.lean}} & first concrete non-self scaled descendant on the windowed route \\
\shortstack[l]{\texttt{Mettapedia/FluidDynamics/}\\\texttt{NavierStokes/WindowedColeHopfHeatSelfCompatibilityObstruction.lean}} & rigidity of same-route closure for the scaled windowed descendant \\
\shortstack[l]{\texttt{Mettapedia/FluidDynamics/}\\\texttt{NavierStokes/WindowedColeHopfHeatSeedBlendFrontier.lean}} & concrete affine seed-blend descendant bridge on the windowed route \\
\shortstack[l]{\texttt{Mettapedia/FluidDynamics/}\\\texttt{NavierStokes/WindowedColeHopfHeatSeedBlendObstruction.lean}} & coefficient rigidity and nonexistence theorems for same-route affine descendants \\
\shortstack[l]{\texttt{Mettapedia/FluidDynamics/}\\\texttt{NavierStokes/WindowedColeHopfHeatFrozenGateFrontier.lean}} & concrete frozen-seed-gated descendant bridge on the windowed route \\
\shortstack[l]{\texttt{Mettapedia/FluidDynamics/}\\\texttt{NavierStokes/WindowedColeHopfHeatFrozenGateObstruction.lean}} & frozen-gate same-route rigidity: nonzero live witnesses force \(a>1\) \\
\shortstack[l]{\texttt{Mettapedia/FluidDynamics/}\\\texttt{NavierStokes/WindowedColeHopfHeatSeedLiveProfileFrontier.lean}} & generic pointwise seed/live-profile shell on the concrete windowed route \\
\shortstack[l]{\texttt{Mettapedia/FluidDynamics/}\\\texttt{NavierStokes/WindowedColeHopfHeatSeedLiveProfileObstruction.lean}} & generic witness law \(P(s,\ell)=\ell\) for same-route pointwise descendants \\
\shortstack[l]{\texttt{Mettapedia/FluidDynamics/}\\\texttt{NavierStokes/WindowedColeHopfHeatSeedLiveOperatorFrontier.lean}} & generic non-pointwise seed/live-operator shell on the concrete windowed route \\
\shortstack[l]{\texttt{Mettapedia/FluidDynamics/}\\\texttt{NavierStokes/WindowedColeHopfHeatSeedLiveOperatorObstruction.lean}} & generic witness law \(O(s,\ell,x)=\ell(x)\) for same-route operator descendants \\
\shortstack[l]{\texttt{Mettapedia/FluidDynamics/}\\\texttt{NavierStokes/WindowedColeHopfHeatSeedOnlyOperatorObstruction.lean}} & seed-only operator descendants force time-invariant witnesses, hence fail on genuinely time-varying states \\
\shortstack[l]{\texttt{Mettapedia/FluidDynamics/}\\\texttt{NavierStokes/WindowedColeHopfHeatLiveOnlySampleKernelObstruction.lean}} & live-only sampled and shift kernels fail when all sampled live witnesses vanish but the target live value does not \\
\shortstack[l]{\texttt{Mettapedia/FluidDynamics/}\\\texttt{NavierStokes/WindowedColeHopfHeatSampleKernelFrontier.lean}} & first finite non-pointwise sampled-kernel descendant on the windowed route \\
\shortstack[l]{\texttt{Mettapedia/FluidDynamics/}\\\texttt{NavierStokes/WindowedColeHopfHeatShiftKernelFrontier.lean}} & first concrete spatial shift-kernel descendant on the windowed route \\
\shortstack[l]{\texttt{Mettapedia/FluidDynamics/}\\\texttt{NavierStokes/WindowedColeHopfHeatShiftKernelObstruction.lean}} & witness-forced same-route coefficient rigidity for the anchored spatial descendant \\
\shortstack[l]{\texttt{Mettapedia/FluidDynamics/}\\\texttt{NavierStokes/WindowedColeHopfHeatMaskedShiftKernelObstruction.lean}} & finite masked spatial families collapse back to the anchored shift obstruction, with bridge-level impossibility corollaries \\
\shortstack[l]{\texttt{Mettapedia/FluidDynamics/}\\\texttt{NavierStokes/WindowedColeHopfHeatSampleKernelObstruction.lean}} & diagonal sampled-kernel same-route rigidity inherited inside the first finite family \\
\shortstack[l]{\texttt{Mettapedia/FluidDynamics/}\\\texttt{NavierStokes/WindowedColeHopfHeatIdentitySampleKernelObstruction.lean}} & any identity-anchor finite sampled kernel collapses back to affine same-route rigidity \\
\shortstack[l]{\texttt{Mettapedia/FluidDynamics/}\\\texttt{NavierStokes/WindowedColeHopfHeatSaturatedFrontier.lean}} & concrete nonlinear saturated descendant bridge on the windowed route \\
\shortstack[l]{\texttt{Mettapedia/FluidDynamics/}\\\texttt{NavierStokes/WindowedColeHopfHeatSaturatedObstruction.lean}} & first nonlinear same-route rigidity theorem on the windowed route \\
\shortstack[l]{\texttt{Mettapedia/FluidDynamics/}\\\texttt{NavierStokes/WindowedColeHopfHeatSaturatedMaskedShiftKernelObstruction.lean}} & widened finite spatial masked-shift shell inherits the nonlinear saturated rigidity \\
\shortstack[l]{\texttt{Mettapedia/FluidDynamics/}\\\texttt{NavierStokes/WindowedColeHopfHeatSaturatedIdentitySampleKernelObstruction.lean}} & identity-anchor finite sampled kernels inherit the nonlinear saturated constant-magnitude obstruction \\
\shortstack[l]{\texttt{Mettapedia/FluidDynamics/}\\\texttt{NavierStokes/WindowedColeHopfHeatSaturatedShiftKernelObstruction.lean}} & arbitrary finite translation kernels inherit the nonlinear saturated constant-magnitude obstruction \\
\shortstack[l]{\texttt{Mettapedia/FluidDynamics/}\\\texttt{NavierStokes/WindowedColeHopfHeatSaturatedShiftKernelMaskedSpecialization.lean}} & the older masked common-shift shell is now derived back from the general nonlinear finite shift-kernel theorem \\
\shortstack[l]{\texttt{Mettapedia/FluidDynamics/}\\\texttt{NavierStokes/FiniteModeLimit.lean}} & finite-mode positivity / energy transfer to the limit \\
\shortstack[l]{\texttt{Mettapedia/FluidDynamics/}\\\texttt{NavierStokes/FiniteModeKernelLimit.lean}} & finite-mode limit transfer into the local kernel shell \\
\shortstack[l]{\texttt{Mettapedia/FluidDynamics/}\\\texttt{NavierStokes/NavierStokesFiniteModeGalerkinToy.lean}} & finite-dimensional Galerkin ODE audit and simple solution families \\
\shortstack[l]{\texttt{Mettapedia/FluidDynamics/}\\\texttt{NavierStokes/ManuscriptTruncationPackage.lean}} & combined manuscript-shaped truncation target \\
\shortstack[l]{\texttt{Mettapedia/FluidDynamics/}\\\texttt{NavierStokes/NavierStokesWitnessConstruction.lean}} & concrete finite-time audit theorem for time-independent compatibility seeds \\
\texttt{ns\_crux.tex} & companion blocker note, kept separate on purpose \\
\bottomrule
\end{tabular}
\end{center}

So the honest answer to ``how much can we build on Mettapedia?'' is:
\begin{quote}
quite a lot for the finite algebra and interface design, not much yet for the
heavy PDE/stochastic geometry itself.
\end{quote}

The current NS-specific Lean files are deliberately modest.  They do not attempt
Kunita theory, infinite-dimensional geometry, or PDE well-posedness.  They
extract the clean finite algebra, the exact hypothesis interface that the
manuscript's identity-only argument depends on, and now a local kernel shell
for future semigroup-facing work.

\section{The New Lean Kernel}

The new module
\texttt{Mettapedia/FluidDynamics/NavierStokes/IdentityEnergy.lean}
defines a finite carré-du-champ style energy
\[
  \Gamma(D) := \sum_i D_i^2
\]
for a finite family of directional derivatives or coefficients.

\subsection{Log-Energy Reduction}

The first kernel theorem is the exact positivity estimate used in the manuscript
when passing from \(\Phi\) to \(\log \Phi\).

\begin{theorem}[Finite log-gradient bound]
If \(m > 0\) and \(m \le F\), then
\[
  \Gamma(D/F) \le \Gamma(D) / m^2.
\]
\end{theorem}

\noindent
This is formalized as \texttt{gamma\_div\_le}.  It is elementary, but it
extracts one of the manuscript's recurrent steps into a reusable statement with
no stochastic or geometric baggage.

The second theorem packages the specific Cole--Hopf scaling used in the NS
route.

\begin{theorem}[The \(-2\nu \log\) energy factor]
If \(m > 0\) and \(m \le \Phi\), then
\[
  \Gamma\!\left((-2\nu)\,D/\Phi\right)
  \le \frac{4\nu^2}{m^2}\,\Gamma(D).
\]
\end{theorem}

\noindent
This is \texttt{gamma\_negTwoNuLog\_le}.  It is exactly the algebraic part of
the passage from a bound on \(\Gamma(P_t \phi)(\mathrm{id})\) to a bound on
\(\Gamma(W(t))(\mathrm{id})\) once \(W(t) = -2\nu \log(P_t \phi)\) and a
positive lower bound for \(P_t \phi\) are available.

\subsection{Positive Lower-Bound Gate}

The companion module
\texttt{Mettapedia/FluidDynamics/NavierStokes/ColeHopfPositiveLowerBound.lean}
turns that phrase ``positive lower bound'' into a reusable Lean object
\texttt{HasColeHopfPointwiseLowerBound}.  It proves the forward consequences
needed by the log layer:
\[
  0 < m \wedge \forall x,\; m \le \Phi(x)
  \quad\Longrightarrow\quad
  0 < \Phi(x), \quad \Phi(x) \ne 0,
\]
and exposes \texttt{gamma\_div\_le} and
\texttt{gamma\_negTwoNuLog\_le} as methods on that certificate.

The same file also proves two converse obstructions relevant to repairs of the
manuscript route.  First,
\[
  \Phi(x_0) \le 0
  \quad\Longrightarrow\quad
  \neg \exists m,\; \texttt{HasColeHopfPointwiseLowerBound}\;\Phi\;m.
\]
In particular, a zero value already rules out the lower-bound certificate.
Second, Lean defines \texttt{HasArbitrarilySmallColeHopfValues} and proves
\[
  \bigl(\forall m>0,\; \exists x,\; \Phi(x)<m\bigr)
  \quad\Longrightarrow\quad
  \neg \exists m,\; \texttt{HasColeHopfPointwiseLowerBound}\;\Phi\;m.
\]
So pointwise positivity alone is still insufficient when the field is not
bounded away from zero.

Lean packages this back into the existing identity interface as
\texttt{ColeHopfIdentityData.Phi\_pos},
\texttt{ColeHopfIdentityData.Phi\_ne\_zero}, the zero-valued route-facing
theorem \texttt{not\_exists\_ColeHopfIdentityData\_of\_zero\_phi}, and the
stronger theorem \texttt{not\_exists\_ColeHopfIdentityData\_of\_arbitrarilySmall\_phi}.
The regression file includes a concrete positive-ray model
\(\Phi(x)=x\) on \(\{x:\mathbb{R}\mid x>0\}\): every value is positive, but
values are arbitrarily small, so no identity-level Cole--Hopf package can use
it.

The latest quotient-level pass adds
\texttt{HasUniformAbsDerivativeLowerBound}.  It proves that if \(\Phi\) is
pointwise positive, has arbitrarily small values, and a scalar derivative
\(d\Phi\) satisfies \(|d\Phi|\ge a>0\), then for every real \(M>0\) there is
some point where
\[
  M < |d\Phi(x)/\Phi(x)|.
\]
Consequently there is no finite uniform bound on this log derivative.  The same
argument is lifted to the one-direction finite \(\gamma\)-energy shell:
\texttt{not\_exists\_uniform\_singleton\_logDerivative\_gamma\_bound\_of\_arbitrarilySmallValues}
rules out a uniform real bound on the squared log-derivative energy.  The
positive-ray regression instantiates the derivative lower bound with derivative
\(1\), so the blocker is not only an interface mismatch; it is a concrete
denominator blow-up model.

The same Lean file now records the exact proportional-vanishing obligation on
the other side.  If \(\Phi(x)>0\), \(\Phi\) has arbitrarily small values, and
there is a uniform bound
\[
  |d\Phi(x)/\Phi(x)| \le B,
\]
then \texttt{abs\_derivative\_arbitrarilySmall\_of\_uniform\_abs\_logDerivative\_bound}
proves that \(d\Phi\) has arbitrarily small absolute values.  Hence
\texttt{not\_exists\_absDerivativeLowerBound\_of\_uniform\_abs\_logDerivative\_bound}
rules out every positive lower bound on \(|d\Phi|\).  The regression suite
includes the proportional positive-ray model \(d\Phi=\Phi=x\), where the
quotient is uniformly \(1\) and the derivative necessarily tends to zero with
\(\Phi\).  This does not solve the NS route; it isolates the minimal analytic
burden that any such repair must discharge for the manuscript's actual heat
field.

\subsection{Frame Push-Down and Vorticity}

The second half of the file formalizes the finite Cauchy--Schwarz step behind
the vorticity estimate.  Given coefficients \(a_i\) and a frame \(E_i(x)\), define
\[
  \mathrm{frameEval}(a,E,x) := \sum_i a_i E_i(x).
\]

\begin{theorem}[Finite frame Cauchy--Schwarz]
For every \(x\),
\[
  |\mathrm{frameEval}(a,E,x)|
  \le \sqrt{\Gamma(a)}\;\sqrt{\Gamma(E(\cdot,x))}.
\]
\end{theorem}

\noindent
This appears in Lean as \texttt{abs\_frameEval\_le}.  The specialized theorem
\texttt{abs\_vorticity\_le} is the manuscript-shaped version where the frame
components are \(\mathrm{curl}\,E_i(x)\).

So the current Lean kernel already isolates the exact positive shape:
\begin{quote}
identity coefficient energy + uniform frame-curl energy
\(\Longrightarrow\) pointwise vorticity bound.
\end{quote}

That is not the whole Navier--Stokes proof.  It is the algebraic spine that the
hard analytic work has to feed.

\subsection{The Cole--Hopf Identity Interface}

The second new module,
\texttt{Mettapedia/FluidDynamics/NavierStokes/ColeHopfInterface.lean},
packages the manuscript-shaped hypotheses into one structure:
\begin{itemize}[leftmargin=1.5em]
  \item a positive profile \(\Phi(t)\),
  \item directional derivatives \(D_i\Phi(t)\) at the identity,
  \item a uniform lower bound \(m_\phi\),
  \item a uniform identity-energy bound on \(\Gamma(D\Phi(t))\),
  \item a curl frame with a uniform frame-energy bound.
\end{itemize}

From that interface, Lean now derives the full algebraic bridge:
\begin{enumerate}[leftmargin=1.5em]
  \item \texttt{gamma\_Wcoeff\_le}: a uniform bound on the coefficients of
  \(W(t) = -2\nu\log\Phi(t)\);
  \item \texttt{abs\_vorticity\_le}: a pointwise vorticity bound at any time and
  point;
  \item \texttt{abs\_vorticity\_le\_uniform}: the uniform version over all times
  and points.
\end{enumerate}

This is exactly the kind of reuse Mettapedia is already good at: small abstract
interfaces with derived consequences, leaving the deep analytic hypotheses
explicit rather than implicit.

\subsection{A Local Kernel-Founded Semigroup Shell}

The next new module,
\texttt{Mettapedia/FluidDynamics/NavierStokes/KernelSemigroupInterface.lean},
does not try to prove SG diffusion facts.  It does something more conservative:
it re-expresses the NS evolution layer in terms of Mettapedia's own stable
kernel semantics from
\texttt{Mettapedia/ProbabilityTheory/MarkovCategory/Kernels.lean}.

Concretely, the file introduces a minimal local structure:
\begin{itemize}[leftmargin=1.5em]
  \item a measurable state object \(\Omega\),
  \item time-indexed endokernels \(P_t : \Omega \rightsquigarrow \Omega\),
  \item identity and composition laws for the kernel family,
  \item the same identity-only Cole--Hopf fields already used in the algebraic
  interface.
\end{itemize}

Lean then proves that all existing coefficient and vorticity consequences are
recovered unchanged from this kernel-founded package.  So the current NS lane
is now grounded on a local semigroup shell without pretending that the hard
analytic semigroup estimates are already done.

\subsection{A Generic Cutoff Continuity Repair Shell}

The newest local helper,
\texttt{Mettapedia/FluidDynamics/NavierStokes/CutoffContinuity.lean},
extracts the exact mechanism behind the chart-cutoff issue without committing
to any SG-specific geometry.

It defines a generic cutoff potential
\[
  W(x) = \eta(r(x))\,s(x)
\]
and proves two small but useful facts:
\begin{enumerate}[leftmargin=1.5em]
  \item if the cutoff, radius observable, and statistic are continuous, then
  \(W\) is continuous;
  \item if a convergent sequence has the cutoff factor tending to the wrong
  limit, then \(W\) cannot be continuous at the limit point.
\end{enumerate}

This helps the route constructively.  The grassroots side can now say:
\begin{quote}
any repair of the chart-cutoff step must supply continuity of the actual radius
observable in the chosen topology, or replace the cutoff construction entirely.
\end{quote}

That does not solve the crux.  It does turn the crux into a local reusable
interface theorem instead of leaving it as manuscript-level warning prose.

\subsection[A Local W-to-phi Topology Shell]{A Local \(W \mapsto \phi\) Topology Shell}

The next new module,
\texttt{Mettapedia/FluidDynamics/NavierStokes/ColeHopfTopology.lean},
pushes the same idea one step further.  It formalizes the real Cole--Hopf map
\[
  \phi = \exp\!\left(-\frac{W}{2\nu}\right)
\]
as a local topological wrapper over the generic cutoff potential.

Its role is again deliberately modest:
\begin{enumerate}[leftmargin=1.5em]
  \item prove that if \(W\) is continuous, then the induced \(\phi\) is
  continuous;
  \item prove that if a sequence forces the wrong limit for \(W\), then it also
  forces the wrong limit for \(\phi\), provided \(\nu \neq 0\);
  \item specialize this to the chart-cutoff potential
  \(W(x)=\eta(r(x))\,s(x)\).
\end{enumerate}

This matters because the crux note's instability is not only about \(W_0\).  It
propagates through the manuscript's actual Cole--Hopf transform.  The new file
turns that propagation into a local reusable theorem rather than leaving it as
an informal remark.

\subsection{A Positive Approximation Repair Shell}

The newest module,
\texttt{Mettapedia/FluidDynamics/NavierStokes/ApproximationInterface.lean},
records the constructive side of the same issue.

It introduces a tiny structure
\texttt{CutoffApproximationData(radius, statistic)}
whose content is exactly this:
\begin{itemize}[leftmargin=1.5em]
  \item an approximant family \(x \mapsto x_n\),
  \item convergence of the radius observable \(r(x_n) \to r(x)\),
  \item convergence of the statistic \(s(x_n) \to s(x)\).
\end{itemize}

From those hypotheses Lean now derives:
\begin{enumerate}[leftmargin=1.5em]
  \item convergence of the cutoff factor \(\eta(r(x_n)) \to \eta(r(x))\);
  \item convergence of the cutoff potential
  \(W(x_n)=\eta(r(x_n))\,s(x_n)\to W(x)\);
  \item convergence of the corresponding Cole--Hopf field
  \(\phi(x_n) \to \phi(x)\);
  \item inheritance of any eventual absolute truncation bound for \(W(x_n)\) or
  \(\phi(x_n)\) by the limit object.
\end{enumerate}

This is the clean positive counterpart to the crux note's mismatch argument.
It says the route does not need vague reassurance about truncations.  It needs
a theorem of one specific shape:
\begin{quote}
the manuscript's actual finite-mode approximants must force convergence of the
radius and statistic observables in the topology being used.
\end{quote}

If that theorem exists, the cutoff/Cole--Hopf passage is no longer handwavy.
The new Lean shell is where such a repair theorem plugs in.  It also means that
``finite-mode first, infinite-mode later'' is no longer just a strategy slogan:
once the finite truncations have a uniform bound and the relevant observables
converge, the same bound transfers to the limit inside Lean.

\subsection{A Concrete D.4 Chart Datum}

The new module
\texttt{Mettapedia/FluidDynamics/NavierStokes/ManuscriptChartDatum.lean}
turns the manuscript's D.4 formula shape into an actual finite-mode Lean model
instead of leaving it only as an interface.

It starts from the existing geometric mode state and defines:
\begin{enumerate}[leftmargin=1.5em]
  \item a squared weighted chart radius \(r^2\), closer to the manuscript's
  \(\|\kappa(g)\|_{H^m}^2\);
  \item a bounded linear observable \(\ell\) with explicit envelope control;
  \item the concrete chart datum \(W_0 = \chi(r^2)\,\ell\);
  \item the concrete heat datum
  \(\phi = \exp(-W_0/(2\nu))\).
\end{enumerate}

Lean then proves three kinds of facts.
\begin{enumerate}[leftmargin=1.5em]
  \item Finite-mode truncation convergence: the squared radius, the bounded
  observable, \(W_0\), and \(\phi\) all converge along truncations.
  \item Boundedness of the observable and the chart datum: once the cutoff
  \(\chi\) is bounded, the resulting \(W_0\) is bounded explicitly.
  \item Lower and upper bounds for \(\phi\): the local D.4 positivity envelope
  is now derived inside Lean from the \(W_0\) bound.
\end{enumerate}

This is still not the manuscript's full SG chart.  But it is a meaningful
upgrade in fidelity.  The chart side is no longer only ``some abstract radius
and statistic.''  There is now a concrete local model where the D.4
construction is correct, bounded, and truncation-stable.

The follow-up file
\texttt{Mettapedia/FluidDynamics/NavierStokes/ManuscriptChartPackage.lean}
then plugs this concrete D.4 chart model straight into the already-existing
\texttt{ManuscriptTruncationPackage.lean}.  So the chart side is not only
concrete in isolation; it now lands directly in the combined package that also
carries the kernel-founded vorticity consequences.

The next new file,
\texttt{Mettapedia/FluidDynamics/NavierStokes/SharedApproximationPackage.lean},
then removes the last interface duplication.  It records one shared
\(\mathrm{approx} : \mathbb{N} \to G \to G\) and derives from that same family:
\begin{enumerate}[leftmargin=1.5em]
  \item the chart-side cutoff/Cole--Hopf approximation package;
  \item the finite-mode identity-energy package;
  \item the combined manuscript-shaped truncation package.
\end{enumerate}

So the current grassroots architecture is now cleaner:
\begin{quote}
one approximation family, one chart side, one identity side, one combined
package.
\end{quote}

The next file,
\texttt{Mettapedia/FluidDynamics/NavierStokes/GeometricSharedApproximation.lean},
is the first actual instance of that architecture.  It takes the concrete D.4
chart datum from the geometric mode state, uses the same truncation family
\texttt{truncateModes}, reads \(d\Phi\) off finitely many selected coefficients
of the same truncated state, and proves:
\begin{enumerate}[leftmargin=1.5em]
  \item eventual positivity for the same \(\Phi\)-family;
  \item eventual finite \(\Gamma(d\Phi)\) bounds for the same family;
  \item direct inhabitation of the full
  \texttt{SharedApproximationPackage.lean}.
\end{enumerate}

So the current grassroots lane no longer says only
``if someone supplies a shared truncation family, we know what to do.''  It now
says:
\begin{quote}
here is one concrete shared family, and here are the positivity and energy
bounds for that same family.
\end{quote}

The next file,
\texttt{Mettapedia/FluidDynamics/NavierStokes/GeometricColeHopfSharedApproximation.lean},
pushes one step closer to the manuscript.  It replaces the toy positive-offset
\(\Phi\) with the actual chart heat datum
\[
  \phi = \exp\!\left(-\frac{W_0}{2\nu}\right)
\]
from the D.4-style chart model, and it defines \(d\Phi\) as that same heat
datum multiplied by finitely selected coefficients of the same truncated state.

That matters because the positivity and energy bounds are now proved for a
family whose \(\Phi\)-side is genuinely Cole--Hopf-shaped.  Concretely, Lean
now proves for the same shared truncation family:
\begin{enumerate}[leftmargin=1.5em]
  \item eventual lower positivity of \(\phi\) from the explicit \(W_0\) bound;
  \item eventual finite \(\Gamma(d\Phi)\) bounds for the \(\phi\)-scaled
  selected-coefficient model;
  \item direct inhabitation of the full
  \texttt{SharedApproximationPackage.lean} by this more faithful choice.
\end{enumerate}

So the situation is now better than before:
\begin{quote}
we no longer have only a shared truncation toy model; we also have a shared
Cole--Hopf-shaped heat-datum model.
\end{quote}

The follow-up file
\texttt{Mettapedia/FluidDynamics/NavierStokes/GeometricColeHopfPackage.lean}
then turns that instance into named concrete package objects.  It exposes:
\begin{enumerate}[leftmargin=1.5em]
  \item the concrete manuscript-shaped truncation package for the geometric
  Cole--Hopf instance;
  \item the associated concrete finite-mode kernel data and kernel semigroup
  data;
  \item the named vorticity field for that package;
  \item direct theorems for cutoff-potential convergence, Cole--Hopf
  convergence, coefficient-energy control, and uniform vorticity bounds.
\end{enumerate}

That is a real usability step.  Later formalization no longer needs to unpack
the whole shared-approximation tower just to cite the first concrete
Cole--Hopf-shaped NS package.

The newest file,
\texttt{Mettapedia/FluidDynamics/NavierStokes/GeometricColeHopfHeatApproximation.lean},
then takes one more real step toward the manuscript's evolution story without
pretending to have proved the full SG semigroup.  It keeps the same D.4 chart
heat datum \(\phi\), keeps the same shared truncation family, but replaces the
static selected-coefficient identity model by
\[
  d\Phi_i(t) = \phi \, e^{-t\,k_i}\, a_{k_i}
\]
for nonnegative time and selected modes \(k_i\).  Lean proves:
\begin{enumerate}[leftmargin=1.5em]
  \item truncation convergence for the heat-decayed coefficient model;
  \item a uniform finite \(\Gamma(d\Phi(t))\) bound, since the heat factors are
  bounded by \(1\) for \(t \ge 0\);
  \item direct inhabitation of the same
  \texttt{SharedApproximationPackage.lean} by this more genuinely
  time-dependent family.
\end{enumerate}

So the identity side is no longer only
\begin{quote}
static selected coefficients multiplied by \(\phi\).
\end{quote}
It is now
\begin{quote}
selected coefficients with an explicit nonnegative-time heat decay, still tied
to the same concrete chart datum and the same finite-mode truncation family.
\end{quote}

\subsection{A Finite-Mode Identity-Energy Limit Shell}

The next new module,
\texttt{Mettapedia/FluidDynamics/NavierStokes/FiniteModeLimit.lean},
does the same kind of cleanup for the manuscript's identity-energy lane.

It proves three things:
\begin{enumerate}[leftmargin=1.5em]
  \item pointwise convergence of a finite coefficient family implies convergence
  of its finite carré-du-champ energy \(\Gamma\);
  \item eventual finite-mode energy bounds \(\Gamma(D_n)\le A\) pass to the
  limit coefficients;
  \item eventual positivity lower bounds \(m \le \Phi_n\) also pass to the
  limit profile.
\end{enumerate}

Those ingredients are then packaged into a structure
\texttt{FiniteModeColeHopfData} whose theorem
\texttt{toColeHopfIdentityData} produces the exact existing
\texttt{ColeHopfIdentityData} interface on the limit object.  So once a future
analytic argument supplies:
\begin{quote}
finite-mode convergence, eventual lower positivity, and eventual coefficient
energy bounds,
\end{quote}
the already verified vorticity bridge applies to the limit automatically.

This is the manuscript's ``finite-mode first, infinite-mode later'' pattern in
a sharper form than before.  It no longer stops at convergence of observables.
It also transfers the actual identity-energy hypotheses needed by the current
Lean vorticity route.

The follow-up module
\texttt{Mettapedia/FluidDynamics/NavierStokes/FiniteModeKernelLimit.lean}
then combines this limit-transfer shell with the already local
\texttt{KernelSemigroupInterface.lean}.  So if a future analytic argument
supplies both a kernel semigroup and the finite-mode convergence/bound package,
the present kernel-founded vorticity estimate is available on the limit object
with no further algebra.

Finally,
\texttt{Mettapedia/FluidDynamics/NavierStokes/ManuscriptTruncationPackage.lean}
packages the full current target in one place.  It combines:
\begin{enumerate}[leftmargin=1.5em]
  \item chart-observable convergence for the cutoff side;
  \item cutoff-potential and Cole--Hopf convergence consequences;
  \item finite-mode positivity and identity-energy transfer to the limit;
  \item the resulting kernel-founded vorticity bound on the limit data.
\end{enumerate}

That matters because the next analytic theorem no longer has to be described
vaguely as ``something feeding several shells.''  It now has one exact Lean
landing point.

\subsection{A Concrete Windowed Route}

The newest concrete pass tests whether the present positive route can survive a
genuine smooth chart window rather than only monotone saturations.  The new
file \texttt{Mettapedia/FluidDynamics/NavierStokes/WindowedCutoffApproximation.lean}
defines the window
\[
  u \mapsto u \cdot
  \bigl(\mathrm{smoothTransition}(cu+1)-\mathrm{smoothTransition}(cu)\bigr)
\]
and proves that it still passes through the existing approximation shell.  So
leaving the stricter \texttt{EnergyProfile} shell does \emph{not} by itself
break chart-side truncation convergence or the induced Cole--Hopf convergence.

The next file,
\texttt{Mettapedia/FluidDynamics/NavierStokes/WindowedColeHopfHeatPackage.lean},
shows that the same window also survives the current quantitative heat-envelope
shell.  Lean proves the explicit bound
\[
  |\chi_c(u)| \le c^{-1}
\]
for the smooth window cutoff \(\chi_c\), and uses that bound to instantiate the
existing heat-decayed Cole--Hopf package.  So there is now a concrete windowed
shared approximation package and a concrete uniform vorticity tendril on that
route.

The next two files then freeze the actual candidate fields carried by that
route.  In
\texttt{Mettapedia/FluidDynamics/NavierStokes/WindowedColeHopfHeatFoundationFrontier.lean},
velocity is fixed to the package vorticity field and pressure is fixed to zero,
with regularity, dynamics, and energy certificates packaged into an
\texttt{AlmostFeffermanBridge}.  In
\texttt{Mettapedia/FluidDynamics/NavierStokes/WindowedColeHopfHeatSeededModel.lean},
the target is strengthened beyond boundedness-only form: pressure must vanish,
velocity must match its distinguished initial slice at time \(1\), and the
velocity bound is made explicit.  So the current concrete route no longer stops
at a raw vorticity tendril; it already reaches a proof-carrying seeded local
model.

The next two frontiers then build genuine non-self descendants on that same
concrete windowed route.  In
\texttt{Mettapedia/FluidDynamics/NavierStokes/WindowedColeHopfHeatScaledSeededFrontier.lean},
Lean packages the scaled descendant \(u(t,x)=a\,\omega(t,x)\).  In
\texttt{Mettapedia/FluidDynamics/NavierStokes/WindowedColeHopfHeatSeedBlendFrontier.lean},
it packages the affine seed-blend descendant
\[
  u(t,x)=a\,\omega(t,x)+b\,\omega(1,x).
\]
The next file,
\texttt{Mettapedia/FluidDynamics/NavierStokes/WindowedColeHopfHeatFrozenGateFrontier.lean},
adds the first seeded nonlinear modulation on that same concrete route:
\[
  u(t,x)=a\,\frac{|\omega(1,x)|}{1+|\omega(1,x)|}\,\omega(t,x).
\]
So the windowed route now reaches not only scalar and affine descendants, but
also a genuinely seed-gated pointwise nonlinearity.  The follow-up obstruction
file
\texttt{Mettapedia/FluidDynamics/NavierStokes/WindowedColeHopfHeatFrozenGateObstruction.lean}
then proves that original self-compatibility already forces
\[
  a\,\frac{|\omega(1,y)|}{1+|\omega(1,y)|}=1
\]
at every nonzero live witness.  Because the frozen gate is always strictly less
than \(1\), any such witness forces \(a>1\).  So on any genuinely nonzero live
state there can be no same-route frozen-gate bridge of this literal form when
\(a \le 1\).  Lean now also proves that any two nonzero live witnesses on such
a bridge must see the same frozen gate.  So the first seeded nonlinear
pointwise route is already rigid on the seed side as well, not only on the
global coefficient.

These pointwise descendants are now unified in
\texttt{Mettapedia/FluidDynamics/NavierStokes/WindowedColeHopfHeatSeedLiveProfileFrontier.lean}.
That file packages the generic shell
\[
  u(t,x)=P(\omega(1,x),\omega(t,x))
\]
directly on the concrete windowed route and proves that the scaled, affine
seed-blend, saturated, and frozen-gate descendants are all exact special cases
of it.  The companion obstruction file
\texttt{Mettapedia/FluidDynamics/NavierStokes/WindowedColeHopfHeatSeedLiveProfileObstruction.lean}
then isolates the shared same-route witness law: at any nonzero live witness,
original self-compatibility forces
\[
  P(\omega(1,y),\omega(t,y))=\omega(t,y).
\]
So the pointwise repair search is no longer only a bag of examples.  It now has
one concrete generic mouth on the actual windowed package, and every literal
pointwise descendant has to close through that same law.

The next compression step is now formalized too in
\texttt{Mettapedia/FluidDynamics/NavierStokes/WindowedColeHopfHeatSeedLiveOperatorFrontier.lean}.
That file packages the full non-pointwise shell
\[
  u(t,x)=O(\omega(1,\cdot),\omega(t,\cdot),x)
\]
on the same concrete windowed route.  It proves that the pointwise
\texttt{SeedLiveProfile} shell, the finite sampled-kernel shell, and the finite
shift-kernel shell are all exact specializations of this one operator frontier.
The companion obstruction file
\texttt{Mettapedia/FluidDynamics/NavierStokes/WindowedColeHopfHeatSeedLiveOperatorObstruction.lean}
then isolates the shared same-route law:
\[
  O(\omega(1,\cdot),\omega(t,\cdot),y)=\omega(t,y).
\]
So even the non-pointwise repair search is no longer a loose collection of
sampled-kernel variants.  On the concrete windowed package it now has one
generic operator mouth, and the finite sampled-kernel and shift-kernel
descendants all feed into it.

That generic operator mouth is now lowered to the explicit finite-frame
componentwise surface too.  In
\texttt{WindowedColeHopfHeatSeedLiveOperatorFrontier.lean}, the new abbreviations
\texttt{windowedColeHopfHeatSeedLiveOperatorInitialSlice\_of\_componentwise\_abs\_le}
and
\texttt{windowedColeHopfHeatSeedLiveOperatorCandidate\_of\_componentwise\_abs\_le}
package the concrete operator shell directly on a componentwise curl bound.
Then \texttt{WindowedColeHopfHeatSeedLiveOperatorObstruction.lean} adds the
componentwise bridge-level theorems
\texttt{windowedColeHopfHeatSeedLiveOperator\_eq\_live\_of\_topDownBridge\_of\_componentwise\_abs\_le}
and
\texttt{not\_exists\_windowedColeHopfHeatSeedLiveOperatorTopDownBridge\_of\_operator\_ne\_live\_of\_componentwise\_abs\_le}.
So the shared non-pointwise operator obstruction itself no longer needs a
separate abstract \texttt{gamma}-bound gap either: any concrete operator/live
mismatch now blocks the whole generic operator shell directly on the same
componentwise finite-frame surface.

The next file,
\texttt{Mettapedia/FluidDynamics/NavierStokes/WindowedColeHopfHeatSeedOnlyOperatorObstruction.lean},
pushes that generic operator mouth one step closer to the real compatibility
issue.  It proves that if a concrete operator descendant actually ignores the
live slice, then same-route closure forces
\[
  \omega(t_1,y)=\omega(t_2,y)
\]
at every witness \(y\) and every pair of times \(t_1,t_2\).  So any genuinely
time-varying witness already rules out that whole seed-only non-pointwise
family.  In other words, once the operator shell is stated honestly, Lean
already shows that a successful repair must use the current live state in a real
dynamical way, not merely the frozen seeded slice.  The same file then pushes
this directly into the first finite non-pointwise shell: any sampled kernel with
all live coefficients equal to \(0\) is a seed-only operator, so it is ruled out
immediately on any witness where the concrete windowed route changes in time.
Because finite shift kernels are just sampled kernels with translation anchors,
the same theorem also kills the first genuinely spatial seed-only corridor:
any shift-kernel descendant with all live coefficients \(0\) already fails on a
time-varying witness.

That seed-only obstruction file is now lowered to the explicit finite-frame
componentwise surface as well.  In
\texttt{WindowedColeHopfHeatSeedOnlyOperatorObstruction.lean}, the new
theorems
\texttt{windowedColeHopfHeatSeedOnlyOperator\_vorticity\_eq\_of\_topDownBridge\_of\_componentwise\_abs\_le},
\texttt{not\_exists\_windowedColeHopfHeatSeedOnlyOperatorTopDownBridge\_of\_vorticity\_ne\_of\_componentwise\_abs\_le},
\texttt{windowedColeHopfHeatSeedOnlySampleKernel\_vorticity\_eq\_of\_topDownBridge\_of\_componentwise\_abs\_le},
\texttt{not\_exists\_windowedColeHopfHeatSeedOnlySampleKernelTopDownBridge\_of\_vorticity\_ne\_of\_componentwise\_abs\_le},
\texttt{windowedColeHopfHeatSeedOnlyShiftKernel\_vorticity\_eq\_of\_topDownBridge\_of\_componentwise\_abs\_le},
and
\texttt{not\_exists\_windowedColeHopfHeatSeedOnlyShiftKernelTopDownBridge\_of\_vorticity\_ne\_of\_componentwise\_abs\_le}
show that the operator, sampled-kernel, and shifted-kernel seed-only
corridors no longer need a separate abstract \texttt{gamma}-bound either.  So
any time-varying concrete witness now blocks those seed-only descendants
directly on the same componentwise finite-frame surface.

The next file,
\texttt{Mettapedia/FluidDynamics/NavierStokes/WindowedColeHopfHeatLiveOnlySampleKernelObstruction.lean},
pushes the dual live-dependent finite-kernel obstruction.  If a sampled kernel
ignores the seeded slice entirely, then same-route closure becomes the pure
live-sample equation
\[
  \sum_i a_i\,\omega(t,\alpha_i(y))=\omega(t,y).
\]
So if every sampled live witness vanishes while the target live value does not,
that literal live-only sampled-kernel bridge is impossible.  Because shift
kernels are only sampled kernels with translation anchors, the same argument
immediately rules out the first genuinely spatial live-only shell too.

That live-only obstruction file is now lowered to the explicit finite-frame
componentwise surface as well.  In
\texttt{WindowedColeHopfHeatLiveOnlySampleKernelObstruction.lean}, the new
theorems
\texttt{windowedColeHopfHeatLiveOnlySampleKernel\_eq\_live\_of\_topDownBridge\_of\_componentwise\_abs\_le},
\texttt{windowedColeHopfHeatLiveOnlySampleKernel\_vorticity\_eq\_zero\_of\_topDownBridge\_of\_zero\_sampled\_live\_of\_componentwise\_abs\_le},
\texttt{not\_exists\_windowedColeHopfHeatLiveOnlySampleKernelTopDownBridge\_of\_zero\_sampled\_live\_nonzero\_live\_of\_componentwise\_abs\_le},
\texttt{windowedColeHopfHeatLiveOnlyShiftKernel\_vorticity\_eq\_zero\_of\_topDownBridge\_of\_zero\_shifted\_live\_of\_componentwise\_abs\_le},
and
\texttt{not\_exists\_windowedColeHopfHeatLiveOnlyShiftKernelTopDownBridge\_of\_zero\_shifted\_live\_nonzero\_live\_of\_componentwise\_abs\_le}
show that this first genuinely live-dependent finite-kernel obstruction no
longer needs a separate abstract \texttt{gamma}-bound either.  So both the
sampled and shifted live-only shells now fail directly on the same concrete
componentwise finite-frame surface.

The route now also reaches a first finite non-pointwise family in
\texttt{Mettapedia/FluidDynamics/NavierStokes/WindowedColeHopfHeatSampleKernelFrontier.lean},
where a sampled kernel mixes finitely many seeded and live samples through
anchor maps.  The diagonal one-point kernel is proved to recover the old affine
seed-blend descendant exactly, so this broader family extends the concrete
windowed route without breaking contact with the earlier rigidity theorems.
The new file
\texttt{Mettapedia/FluidDynamics/NavierStokes/WindowedColeHopfHeatSampleKernelObstruction.lean}
pushes that one step further: on the diagonal sampled-kernel case, the affine
same-route bridge obstructions are inherited verbatim.  So a nonzero initial
slice still forces \(a+b=1\), and the stronger witness patterns still force
\(a=1\) or \(b=0\), now inside the first finite non-pointwise family rather
than only in the pointwise affine shell.  The next file,
\texttt{Mettapedia/FluidDynamics/NavierStokes/WindowedColeHopfHeatIdentitySampleKernelObstruction.lean},
shows that this is not merely a one-point diagonal curiosity: whenever a finite
sampled kernel keeps both seed and live anchors equal to the identity, the
whole descendant collapses back to the affine seed-blend route with coefficients
given by the total live and seed masses.  So finite repetition at the same
point still forces the same affine same-route identities on those total masses.
The newest nonlinear step adds
\texttt{Mettapedia/FluidDynamics/NavierStokes/WindowedColeHopfHeatSaturatedFrontier.lean},
which packages the saturated descendant
\[
  u(t,x)=a\,\frac{\omega(t,x)}{1+|\omega(t,x)|}.
\]
All four routes now carry full local certificates and descend to explicit
bridge objects on the concrete windowed package.

The spatial route is now concrete too.  In
\texttt{WindowedColeHopfHeatShiftKernelFrontier.lean}, Lean proves that the
anchored descendant
\[
  u(t,x)=a\,\omega(t,x)+b\,\omega(1,x+s)+d\,\omega(t,x+s)
\]
is exactly a finite translation-kernel instance of the generic shift-kernel
frontier.  In
\texttt{WindowedColeHopfHeatShiftKernelObstruction.lean}, Lean then proves the
first same-route rigidity laws for that spatial descendant: if the shifted
seeded and shifted live slices vanish while the unshifted live slice is
nonzero, then \(a=1\); if the unshifted live and shifted live slices vanish
while the shifted seeded slice is nonzero, then \(b=0\); and if the unshifted
live and shifted seeded slices vanish while the shifted live slice is nonzero,
then \(d=0\).  So the first genuinely spatial descendant is not only packaged,
but already constrained by theorem-level same-route rigidity on the concrete
windowed route.

The next widening step,
\texttt{WindowedColeHopfHeatMaskedShiftKernelObstruction.lean}, lifts that
spatial rigidity from one anchored descendant to a finite Boolean-masked family.
Every seed term uses the same shift \(s\), while each live term either stays at
\(x\) or moves to \(x+s\).  Lean proves that this whole family collapses back to
the same three-coefficient anchored route with coefficients given by the total
unshifted-live mass, total seed mass, and total shifted-live mass.  So the same
witness patterns now force, at finite-family level,
\[
  a_{\mathrm{base}}=1,\qquad b_{\mathrm{seed}}=0,\qquad
  d_{\mathrm{shift}}=0,
\]
and the matching bridge-level impossibility corollaries are formalized too.  So
the first concrete spatial rigidity result is no longer a one-instance fact; it
already controls a nontrivial finite spatial shell on the concrete windowed
route.

Finally, the corresponding obstruction files record that same-route closure is
already rigid on this concrete route.  In
\texttt{WindowedColeHopfHeatSelfCompatibilityObstruction.lean}, Lean proves that
on any genuinely nonzero vorticity state the scaled descendant can satisfy the
original self-compatibility mouth only when \(a=1\).  In
\texttt{WindowedColeHopfHeatSeedBlendObstruction.lean}, Lean proves that
same-route closure for the affine seed-blend descendant forces explicit
coefficient identities: \(a+b=1\) under a nonzero initial slice, and under
stronger witnesses even \(a=1\) or \(b=0\).  So the grassroots route now
contains not only positive concrete descendants, but also theorem-level
rigidity about when those descendants can still pretend to be same-route.  The
same affine mouth is now lowered to the explicit componentwise finite-frame
surface too.  In
\texttt{WindowedColeHopfHeatSeedBlendFrontier.lean}, the new abbreviations
\texttt{windowedColeHopfHeatSeedBlendInitialSlice\_of\_componentwise\_abs\_le}
and
\texttt{windowedColeHopfHeatSeedBlendCandidate\_of\_componentwise\_abs\_le}
package the affine initial slice and candidate directly from a componentwise
curl bound.  In
\texttt{WindowedColeHopfHeatSeedBlendObstruction.lean}, the componentwise
counterparts
\texttt{windowedColeHopfHeatSeedBlend\_sum\_eq\_one\_of\_topDownBridge\_of\_nonzero\_initial\_of\_componentwise\_abs\_le},
\texttt{windowedColeHopfHeatSeedBlend\_a\_eq\_one\_of\_topDownBridge\_of\_zero\_initial\_nonzero\_live\_of\_componentwise\_abs\_le},
\texttt{windowedColeHopfHeatSeedBlend\_b\_eq\_zero\_of\_topDownBridge\_of\_zero\_live\_nonzero\_initial\_of\_componentwise\_abs\_le}
and their nonexistence corollaries lower those first affine same-route
rigidity laws to the same concrete surface.  So this seed-blend obstruction
layer no longer carries a separate abstract \texttt{gamma}-bound gap either.
The
new file
\texttt{WindowedColeHopfHeatSaturatedObstruction.lean} now adds the first
nonlinear rigidity law on the same concrete route: if the saturated descendant
is forced back through original self-compatibility at a nonzero vorticity
witness, then \(a = 1 + |\omega(t,x)|\), hence in particular no such same-route
bridge can exist on a nonzero-vorticity state when \(a \le 1\).  The new
bridge-level upgrade is stronger: any same-route saturated bridge now forces
all nonzero witnesses to have the same absolute vorticity magnitude.  So two
distinct nonzero magnitudes already rule out that nonlinear closure entirely.
The new file
\texttt{WindowedColeHopfHeatSaturatedIdentitySampleKernelObstruction.lean}
pushes that same nonlinear rigidity through every finite sampled kernel whose
seed and live anchors are both the identity.  In that regime the sampled-kernel
candidate still collapses to a total-mass affine combination at time \(1\), so
if it realizes the saturated descendant with \(a \neq 0\), then all nonzero
initial witnesses must again have the same absolute vorticity magnitude.  So
finite repetition at the same point does not rescue the nonlinear saturated
route either.

The corresponding widened finite spatial shell is now pinned down too.  In
\texttt{WindowedColeHopfHeatSaturatedMaskedShiftKernelObstruction.lean}, Lean
proves that if the saturated descendant is realized through the finite masked
shift-kernel family and one has a witness with zero shifted initial slice, zero
shifted live slice, and nonzero unshifted live slice, then the total
unshifted-live mass must satisfy
\[
  a_{\mathrm{base}} \bigl(1 + |\omega(t,x)|\bigr) = a.
\]
So if there are two such witnesses with distinct nonzero absolute vorticity
magnitudes, the widened masked-shift spatial shell cannot realize the saturated
descendant at all.  The same file now also reaches the explicit componentwise
frame-bound surface via the componentwise counterparts
\texttt{windowedColeHopfHeatMaskedShiftKernel\_liveBaseCoeffSum\_mul\_one\_add\_abs\_eq\_saturatedCoeff\_of\_topDownBridge\_of\_zero\_shifted\_initial\_zero\_shifted\_live\_nonzero\_live\_of\_componentwise\_abs\_le},
\texttt{windowedColeHopfHeatMaskedShiftKernel\_abs\_eq\_abs\_of\_topDownBridge\_of\_two\_zero\_shifted\_witnesses\_of\_componentwise\_abs\_le},
and
\texttt{not\_exists\_windowedColeHopfHeatMaskedShiftKernelTopDownBridge\_eq\_saturatedCandidate\_of\_abs\_ne\_abs\_with\_zero\_shifted\_witnesses\_of\_componentwise\_abs\_le}.
Thus the first nonlinear masked spatial escape corridor is already formally
rigid on that same concrete finite-frame surface too.

The next file
\texttt{WindowedColeHopfHeatSaturatedShiftKernelObstruction.lean} removes the
common-shift restriction altogether.  Lean proves that if the saturated
descendant is realized through an arbitrary finite translation-kernel shell and
one has a witness whose seeded anchors all vanish, whose nonzero live anchors
all vanish, and whose unshifted live value remains nonzero, then the total
zero-shift live mass must satisfy
\[
  a_{0} \bigl(1 + |\omega(t,x)|\bigr) = a.
\]
So if two such witnesses survive with distinct nonzero absolute magnitudes, the
entire finite translation-kernel shell cannot realize the saturated descendant.
At that point the nonlinear rigidity is no longer tied to one concrete spatial
encoding; it already lives at the general finite shift-kernel level.  The same
Lean file now also specializes this back to the first anchored translation
kernel, so even the original concrete spatial descendant does not reopen the
nonlinear route once two such witness magnitudes disagree.

That same zero-seed / nonzero-shift obstruction family now also reaches the
explicit componentwise frame-bound surface.  In the same Lean file, the new
theorems
\texttt{windowedColeHopfHeatShiftKernel\_liveZeroCoeffSum\_mul\_one\_add\_abs\_eq\_saturatedCoeff\_of\_topDownBridge\_of\_zero\_seed\_nonzero\_shift\_witness\_of\_componentwise\_abs\_le},
\texttt{windowedColeHopfHeatShiftKernel\_abs\_eq\_abs\_of\_topDownBridge\_of\_two\_zero\_seed\_nonzero\_shift\_witnesses\_of\_componentwise\_abs\_le},
\texttt{not\_exists\_windowedColeHopfHeatShiftKernelTopDownBridge\_eq\_saturatedCandidate\_of\_abs\_ne\_abs\_with\_zero\_seed\_nonzero\_shift\_witnesses\_of\_componentwise\_abs\_le},
and
\texttt{not\_exists\_windowedColeHopfHeatAnchoredTranslationShiftKernelTopDownBridge\_eq\_saturatedCandidate\_of\_abs\_ne\_abs\_of\_componentwise\_abs\_le}
show that this nonlinear witness obstruction no longer needs a separate
abstract \texttt{gamma}-energy hypothesis either.  So the general
finite-kernel saturated obstruction and its first anchored special case now sit
directly on the same concrete finite-frame bound surface as the invariant and
anchored invariant corridors.  The same file now lowers the separate
zero-other-anchors / nonzero-initial branch too via
\texttt{windowedColeHopfHeatShiftKernel\_zeroAnchorCoeffSum\_mul\_one\_add\_abs\_eq\_saturatedCoeff\_of\_topDownBridge\_of\_zero\_other\_anchors\_nonzero\_initial\_of\_componentwise\_abs\_le},
\texttt{windowedColeHopfHeatShiftKernel\_abs\_eq\_abs\_of\_topDownBridge\_of\_two\_zero\_other\_anchors\_nonzero\_initials\_of\_componentwise\_abs\_le},
and
\texttt{not\_exists\_windowedColeHopfHeatShiftKernelTopDownBridge\_eq\_saturatedCandidate\_of\_abs\_ne\_abs\_with\_zero\_other\_anchors\_nonzero\_initials\_of\_componentwise\_abs\_le}.
So the whole generic saturated shift-kernel obstruction file now sits directly
on that concrete finite-frame bound surface.  The remaining burden now moves to
the masked specialization and masked frontier descendants.

The follow-up file
\texttt{WindowedColeHopfHeatSaturatedShiftKernelMaskedSpecialization.lean} now
closes the loop the other way too: it shows that the older masked common-shift
family is already a special case of the newer general finite shift-kernel
theorem.  So the nonlinear spatial obstruction picture is no longer a stack of
separate examples; the older masked shell now sits inside the newer general
one, with the earlier masked impossibility corollary recovered from the more
abstract finite-kernel theorem.  The same specialization file now also reaches
the explicit componentwise frame-bound surface via
\texttt{windowedColeHopfHeatMaskedShiftKernel\_liveBaseCoeffSum\_mul\_one\_add\_abs\_eq\_saturatedCoeff\_of\_topDownBridge\_via\_shiftKernel\_of\_componentwise\_abs\_le}
and
\texttt{not\_exists\_windowedColeHopfHeatMaskedShiftKernelTopDownBridge\_eq\_saturatedCandidate\_of\_abs\_ne\_abs\_with\_zero\_shifted\_witnesses\_via\_shiftKernel\_of\_componentwise\_abs\_le}.
So the masked common-shift shell is no longer only contained in the newer
finite-kernel theorem abstractly; it is already recovered on that same
concrete finite-frame surface too.

The constructive side of that same masked common-shift corridor is now lowered
as well in
\texttt{WindowedColeHopfHeatSaturatedMaskedShiftKernelFrontier.lean}.  The new
theorems
\texttt{windowedColeHopfHeatMaskedShiftKernelCandidate\_eq\_saturatedCandidate\_of\_stationary\_shiftInvariant\_constantMagnitude\_of\_componentwise\_abs\_le},
\texttt{windowedColeHopfHeat\_realizes\_saturated\_pressure\_seeded\_clause\_via\_maskedShiftKernel\_of\_shiftInvariant\_constantMagnitude\_of\_componentwise\_abs\_le},
and
\texttt{exists\_windowedColeHopfHeatSaturatedMaskedShiftKernelBridge\_of\_stationary\_shiftInvariant\_constantMagnitude\_of\_componentwise\_abs\_le}
show that the older masked common-shift frontier no longer needs a separate
abstract \texttt{gamma}-energy bound either.  So this masked branch now sits on
the same concrete finite-frame surface on both its constructive and obstructive
sides.  The remaining burden moves beyond this masked common-shift lane to the
broader nonlinear descendants.

That broader constructive descendant is now lowered too in
\texttt{WindowedColeHopfHeatSaturatedShiftKernelConstantMagnitudeFrontier.lean}.
The new theorems
\texttt{windowedColeHopfHeatSaturatedConstantMagnitudeShiftKernelInitialSlice\_eq\_saturatedInitialSlice\_of\_componentwise\_abs\_le},
\texttt{windowedColeHopfHeatSaturatedConstantMagnitudeShiftKernelCandidate\_eq\_saturatedCandidate\_of\_componentwise\_abs\_le},
\texttt{windowedColeHopfHeat\_realizes\_saturated\_pressure\_seeded\_clause\_via\_constantMagnitudeShiftKernel\_of\_componentwise\_abs\_le},
and
\texttt{exists\_windowedColeHopfHeatSaturatedConstantMagnitudeShiftKernelBridge\_of\_componentwise\_abs\_le}
show that the first honest constructive survivor inside the full finite
shift-kernel shell no longer needs a separate abstract \texttt{gamma}-energy
bound either.  So the zero-shift live-only constant-magnitude corridor now
also sits directly on that same concrete finite-frame surface.  The remaining
burden moves beyond this constructive survivor to the other broader nonlinear
descendants.

The first sampled constructive descendant on top of that survivor is now
lowered too in
\texttt{WindowedColeHopfHeatSaturatedIdentitySampleKernelFrontier.lean}.  The
new theorems
\texttt{windowedColeHopfHeatIdentitySampleKernelInitialSlice\_eq\_saturatedInitialSlice\_of\_constantMagnitude\_of\_componentwise\_abs\_le},
\texttt{windowedColeHopfHeatIdentitySampleKernelCandidate\_eq\_saturatedCandidate\_of\_stationary\_constantMagnitude\_of\_componentwise\_abs\_le},
\texttt{windowedColeHopfHeat\_realizes\_saturated\_pressure\_seeded\_clause\_via\_identitySampleKernel\_of\_constantMagnitude\_of\_componentwise\_abs\_le},
and
\texttt{exists\_windowedColeHopfHeatSaturatedIdentitySampleKernelBridge\_of\_stationary\_constantMagnitude\_of\_componentwise\_abs\_le}
show that every finite sampled kernel whose seed and live anchors are both the
identity already inherits the same concrete componentwise constant-magnitude
saturated bridge.  So the first nonlinear sampled constructive shell no longer
needs a separate abstract \texttt{gamma}-bound either.

That last remaining one-point sampled descendant is now lowered too in
\texttt{WindowedColeHopfHeatSaturatedDiagonalSampleKernelFrontier.lean}.  The
new theorems
\texttt{windowedColeHopfHeatDiagonalSampleKernelInitialSlice\_eq\_saturatedInitialSlice\_of\_constantMagnitude\_of\_componentwise\_abs\_le},
\texttt{windowedColeHopfHeatDiagonalSampleKernelCandidate\_eq\_saturatedCandidate\_of\_stationary\_constantMagnitude\_of\_componentwise\_abs\_le},
\texttt{windowedColeHopfHeat\_realizes\_saturated\_pressure\_seeded\_clause\_via\_diagonalSampleKernel\_of\_constantMagnitude\_of\_componentwise\_abs\_le},
and
\texttt{exists\_windowedColeHopfHeatSaturatedDiagonalSampleKernelBridge\_of\_stationary\_constantMagnitude\_of\_componentwise\_abs\_le}
show that the older diagonal one-sample shell is now also recovered on that
same explicit finite-frame surface.  So the remaining sampled constructive
burden now lies beyond these identity-anchor descendants in the later
nonlinear shells.

The matching nonlinear obstruction for that same older one-sample shell is now
lowered too in
\texttt{WindowedColeHopfHeatSaturatedSampleKernelObstruction.lean}.  The new
theorems
\texttt{windowedColeHopfHeatDiagonalSampleKernel\_sum\_mul\_one\_add\_abs\_eq\_saturatedCoeff\_of\_topDownBridge\_of\_nonzero\_initial\_of\_componentwise\_abs\_le},
\texttt{windowedColeHopfHeatDiagonalSampleKernel\_abs\_eq\_abs\_of\_topDownBridge\_of\_two\_nonzero\_initial\_of\_componentwise\_abs\_le},
and
\texttt{not\_exists\_windowedColeHopfHeatDiagonalSampleKernelTopDownBridge\_eq\_saturatedCandidate\_of\_abs\_ne\_abs\_of\_componentwise\_abs\_le}
show that the diagonal sampled obstruction now sits on the same explicit
finite-frame surface too: any sampled bridge to the saturated descendant
already forces equal absolute values across every nonzero initial witness
pair, so an unequal-absolute-value witness blocks the route without any
separate abstract \texttt{gamma}-bound.

That broader identity-anchor sampled obstruction is now lowered too in
\texttt{WindowedColeHopfHeatSaturatedIdentitySampleKernelObstruction.lean}.
The new theorems
\texttt{windowedColeHopfHeatIdentitySampleKernel\_sum\_mul\_one\_add\_abs\_eq\_saturatedCoeff\_of\_topDownBridge\_of\_nonzero\_initial\_of\_componentwise\_abs\_le},
\texttt{windowedColeHopfHeatIdentitySampleKernel\_abs\_eq\_abs\_of\_topDownBridge\_of\_two\_nonzero\_initial\_of\_componentwise\_abs\_le},
and
\texttt{not\_exists\_windowedColeHopfHeatIdentitySampleKernelTopDownBridge\_eq\_saturatedCandidate\_of\_abs\_ne\_abs\_of\_componentwise\_abs\_le}
show that the first genuinely broader sampled obstruction shell now lands on
that same concrete finite-frame surface as well.  So the identity-anchor
sampled route now has both its constructive frontier and its saturated
obstruction boundary expressed directly in componentwise finite-frame terms,
with no separate abstract \texttt{gamma}-bound gap left in this file either.

That remaining specialization gap is now lowered too in
\texttt{WindowedColeHopfHeatSaturatedShiftKernelIdentitySampleSpecialization.lean}.
The new theorems
\texttt{windowedColeHopfHeatIdentitySampleKernel\_sum\_mul\_one\_add\_abs\_eq\_saturatedCoeff\_of\_topDownBridge\_via\_saturatedShiftKernel\_of\_componentwise\_abs\_le},
\texttt{windowedColeHopfHeatIdentitySampleKernel\_abs\_eq\_abs\_of\_topDownBridge\_via\_saturatedShiftKernel\_of\_two\_nonzero\_initial\_of\_componentwise\_abs\_le},
and
\texttt{not\_exists\_windowedColeHopfHeatIdentitySampleKernelTopDownBridge\_eq\_saturatedCandidate\_of\_abs\_ne\_abs\_via\_saturatedShiftKernel\_of\_componentwise\_abs\_le}
show that the general saturated shift-kernel obstruction already subsumes the
older identity-anchor sampled obstruction directly on the same explicit
finite-frame componentwise surface.  So this specialization layer no longer
contains its own abstract \texttt{gamma}-bound escape hatch either; the
remaining sampled obstruction burden moves past the specialization step and
into later descendants.

\section{Why This Matters}

The value of this kernel is not that it proves a Millennium problem.  The value
is that it separates the route into two layers.

\begin{enumerate}[leftmargin=1.5em]
  \item \textbf{Reusable algebraic layer.}
  Once one has a lower bound on \(P_t \phi\), an identity-energy bound for
  \(\Gamma(P_t\phi)(\mathrm{id})\), and a frame-curl constant, the rest of the
  vorticity estimate is just these finite inequalities.
  \item \textbf{Hard analytic input layer.}
  The real burden is now explicit: produce those bounds in the actual infinite-
  dimensional SG setting, justify the semigroup/flow passage, and justify the
  push-down theorem analytically.
\end{enumerate}

That split is precisely what a grassroots program should want.  If the SG route
survives, these lemmas will still be needed.  If the SG route mutates, the same
identity-energy interface may survive even then.

\section{Positive Manuscript Infrastructure}

Even on a charitable read, not every part of the manuscript deserves the same
confidence.  But the positive architecture is real enough to build around.

\subsection{Polynomial Fourier Frame}

The manuscript gives a concrete divergence-free Fourier frame with polynomial
orbit weights on \(T^3\).  This is not just inspirational prose.  It is a
specific normalization scheme designed to make the relevant \(\ell^2\)-type
sums finite for \(m \ge 5\).  Grassroots formalization can treat this as a real
mathematical object, separate from the later stochastic and geometric layers.

\subsection{Finite-Mode First, Infinite-Mode Later}

The manuscript repeatedly follows a sensible proof pattern:
\begin{quote}
prove the estimate for finite truncations first, then pass to the full
infinite-mode object.
\end{quote}

That is exactly the right shape for formalization.  It suggests that future Lean
work should not start with the full diffeomorphism-group diffusion.  It should
start with finite coefficient families, finite frame sums, and explicit
truncation interfaces.

\subsection{Schwartz Exhaustion Is Background, Not Center Stage}

On a closer read, the manuscript's periodization and torus-exhaustion appendix
is more substantial than a first skim suggests.  For grassroots purposes, the
important point is structural:
\begin{quote}
the \(\mathbb{R}^3\) step mostly wants uniform control of the same torus-side
constants already needed earlier.
\end{quote}

So the best constructive work remains the torus-side background theory, not an
early jump to the Euclidean limit.

\subsection{Concrete Seed Witness Audit}

The finite-time witness surface also now has a bottom-up audit theorem that is
independent of the SG route.  In
\texttt{NavierStokesWitnessConstruction.lean}, Lean proves
\texttt{timeIndependentVelocity\_nonMomentumFiniteTimeFields}: a smooth,
divergence-free, finite-energy initial datum gives a time-independent seed with
smoothness, incompressibility, initial matching, zero time derivative, and slab
energy control.  The missing field is exactly the momentum equation with a
pressure.

The companion theorem
\texttt{timeIndependentVelocity\_finiteTimeWitness\_iff\_stationaryMomentum}
then states the precise boundary: a fixed pressure turns that seed into an
\texttt{ExplicitFiniteTimeRegularityWitness} on a slab if and only if the
stationary momentum balance holds throughout the slab.  The older boxed
partial-periodization steady-seed classification is now an instance of this
general theorem.  Thus boxed compatibility data are documented as useful
local-existence prerequisites, not as an evolved Navier--Stokes solution.

The negative direction is now also available in pointwise audit form:
\texttt{not\_exists\_timeIndependentVelocity\_finiteTimeWitness\_of\_stationaryMomentum\_failure}
and its boxed specialization show that a single stationary residual failure at
one spacetime point of the slab rules out that exact time-independent
velocity/pressure pair as a finite-time witness.  This is deliberately an
exact-seed exclusion.  It does not preclude a local-existence theorem that
constructs a genuinely time-dependent evolved velocity and pressure from the
same initial datum.

The same audit is also exposed in forward form.  The theorems
\texttt{timeIndependentVelocity\_finiteTimeWitness\_implies\_stationaryMomentum}
and
\texttt{timeIndependentVelocity\_finiteTimeWitness\_implies\_stationaryMomentum\_zero\_time}
say that any exact finite-time use of a time-independent seed must satisfy the
stationary residual throughout the slab, and in particular at \(t=0\) on a
nonnegative slab.  The boxed partial-periodization versions provide the same
necessary-condition checklist for the manuscript-style boxed steady seed.

The zero-velocity gauge case has also been split out as a concrete baseline.
Lean proves
\texttt{zeroVelocity\_finiteTimeWitness\_iff\_spatialPressureGradient\_zero}:
zero velocity with a fixed smooth pressure field is a finite-time witness
exactly when that pressure has zero spatial gradient on the slab.  The positive
instance
\texttt{zeroVelocity\_timeOnlyPressure\_finiteTimeWitness} covers smooth
time-only pressure gauges, while
\texttt{not\_exists\_zeroVelocity\_finiteTimeWitness\_of\_spatialPressureGradient\_failure}
turns any nonzero spatial pressure gradient at one slab point into an exact
no-witness certificate for the zero-velocity route.

The same harmless-gauge separation is now explicit for the manuscript-style
boxed steady seed.  The theorem
\texttt{boxedPartialPeriodizationSteadySeed\_zeroSpatialGradientPressure\_finiteTimeWitness\_iff\_stationaryMomentum\_zeroPressure}
already covers the full smooth zero-spatial-gradient gauge class, and
\texttt{boxedPartialPeriodizationSteadySeed\_timeOnlyPressure\_finiteTimeWitness\_iff\_stationaryMomentum\_zeroPressure}
says in particular that adding a smooth time-only pressure gauge \(p(t)\) does
not change the boxed steady-seed witness criterion at all: the exact finite-time
witness condition is still the same zero-pressure stationary momentum balance on
the slab.  The companion theorem
\texttt{not\_exists\_boxedPartialPeriodizationSteadySeed\_zeroSpatialGradientPressure\_finiteTimeWitness\_of\_stationaryMomentum\_failure}
already rules out every such harmless gauge, while
\texttt{not\_exists\_boxedPartialPeriodizationSteadySeed\_timeOnlyPressure\_finiteTimeWitness\_of\_stationaryMomentum\_failure}
then turns any pointwise failure of that zero-pressure stationary residual into
an exact no-witness certificate even after allowing the manuscript's time-only
specialization.  So the boxed steady-seed route cannot be rescued by harmless
zero-spatial-gradient pressure changes; any repair has to modify the velocity
side or solve the stationary balance itself.  Lean now packages that route-level
obstruction directly as
\texttt{boxedPartialPeriodizationSteadySeed\_timeOnlyPressure\_exhibits\_nonMomentumFiniteTimeFields\_without\_exactWitness}:
the boxed seed already carries the smoothness, divergence, initial-match, and
bounded-energy witness fields with a smooth time-only gauge, yet still fails to
upgrade to an exact witness whenever the zero-pressure stationary residual
breaks at one slab point.

That same fixed-pair obstruction now reaches the exact whole-space output
surface too.  In
\texttt{NavierStokesWitnessConstructionGlobalOutputObstruction.lean}, Lean adds
the exact-pair predicate
\texttt{ExplicitConcreteNavierStokesGlobalOutputWithVelocityPressure} and the
bridge
\texttt{ExplicitConcreteNavierStokesGlobalOutputWithVelocityPressure.implies\_finiteTimeWitness},
which says that any exact global output for a fixed velocity/pressure pair
restricts to an exact finite-time witness on every slab.  Its contrapositive
\texttt{not\_ExplicitConcreteNavierStokesGlobalOutputWithVelocityPressure\_of\_not\_exists\_finiteTimeWitness}
therefore lifts any slabwise exact-witness obstruction to the exact whole-space
surface.  The boxed specialization
\texttt{not\_ExplicitConcreteNavierStokesGlobalOutputWithVelocityPressure\_boxedPartialPeriodizationSteadySeed\_timeOnlyPressure\_of\_stationaryMomentum\_failure}
shows that the manuscript-style boxed steady seed with a smooth time-only
pressure gauge cannot even be an exact global output when the same
zero-pressure stationary residual fails at one slab point, and
\texttt{boxedPartialPeriodizationSteadySeed\_timeOnlyPressure\_exhibits\_nonMomentumGlobalFields\_without\_exactGlobalOutput}
records that the seed still carries every other whole-space output field.

That same obstruction now survives the older uniform-vorticity packaging layer
too.  In
\texttt{NavierStokesWitnessConstructionUniformVorticityObstruction.lean}, the
generic lift
\texttt{not\_exists\_uniformVorticityData\_of\_not\_exists\_finiteTimeWitness}
shows that adding a slab vorticity bound cannot rescue a velocity/pressure pair
that already fails the exact finite-time witness obligations.  Its boxed
specializations
\texttt{not\_exists\_boxedPartialPeriodizationSteadySeed\_timeOnlyPressure\_uniformVorticityData\_of\_stationaryMomentum\_failure}
and
\texttt{boxedPartialPeriodizationSteadySeed\_timeOnlyPressure\_exhibits\_nonMomentumFiniteTimeFields\_without\_uniformVorticityData}
show that the manuscript-style boxed steady seed still cannot enter the older
uniform-vorticity premise when the same zero-pressure stationary residual fails
at one slab point.

That same obstruction now survives the repaired BKM packaging layer too.  In
\texttt{NavierStokesWitnessConstructionBKMObstruction.lean}, the generic lift
\texttt{not\_exists\_BKMData\_of\_not\_exists\_finiteTimeWitness} shows that a
vorticity envelope and integral budget cannot rescue a velocity/pressure pair
that already fails the exact finite-time witness obligations.  Its boxed
specializations
\texttt{not\_exists\_boxedPartialPeriodizationSteadySeed\_timeOnlyPressure\_BKMData\_of\_stationaryMomentum\_failure}
and
\texttt{boxedPartialPeriodizationSteadySeed\_timeOnlyPressure\_exhibits\_nonMomentumFiniteTimeFields\_without\_BKMData}
show that the manuscript-style boxed steady seed still cannot enter the
repaired BKM premise when the same zero-pressure stationary residual fails at
one slab point.

The same file now also records the matching positive boundary for Schwartz
time-independent seeds.  The theorem
\texttt{timeIndependentVelocity\_schwartz\_BKMData\_iff\_stationaryMomentum}
proves that for a Schwartz initial velocity with divergence-free initial slice
and a smooth fixed pressure, BKM-packaged data for the exact
time-independent seed exist if and only if the stationary momentum balance
holds on the slab.  The reverse direction uses the single-profile Schwartz
vorticity envelope internally, so the BKM layer is conservative over the
finite-time witness audit without leaving an extra vorticity hypothesis to be
supplied by hand.

The boxed steady-seed surface now has the same exact classification, not only
one-sided obstructions.  Theorems
\texttt{boxedPartialPeriodizationSteadySeed\_uniformVorticityData\_iff\_stationaryMomentum}
and
\texttt{boxedPartialPeriodizationSteadySeed\_BKMData\_iff\_stationaryMomentum}
prove that, for any fixed smooth pressure, the boxed steady seed has
uniform-vorticity-packaged or BKM-packaged finite-time data exactly when the
boxed stationary momentum balance already holds on the slab.  The constructive
directions use the finite boxed Schwartz periodization profile to supply the
vorticity bound/envelope.  Thus these continuation packages do not repair the
boxed steady seed; they only expose the original stationary momentum equation
as the remaining obligation.  The follow-up time-only pressure theorems
\texttt{boxedPartialPeriodizationSteadySeed\_timeOnlyPressure\_uniformVorticityData\_iff\_stationaryMomentum\_zeroPressure}
and
\texttt{boxedPartialPeriodizationSteadySeed\_timeOnlyPressure\_BKMData\_iff\_stationaryMomentum\_zeroPressure}
make the manuscript-facing gauge case explicit: a smooth time-only pressure can
be carried through the packaged data, but it is equivalent to satisfying the
zero-pressure stationary audit.  The existential zero-gradient variants
\texttt{exists\_boxedPartialPeriodizationSteadySeed\_zeroSpatialGradientPressure\_uniformVorticityData\_iff\_stationaryMomentum\_zeroPressure}
and
\texttt{exists\_boxedPartialPeriodizationSteadySeed\_zeroSpatialGradientPressure\_BKMData\_iff\_stationaryMomentum\_zeroPressure}
close the broader harmless-gauge repair class at the packaged-data layer: even
choosing an arbitrary smooth pressure with zero spatial gradient is equivalent
to satisfying that same zero-pressure stationary audit.

These four exact-pair escape hatches are now collapsed into one certified route
surface.  In
\texttt{NavierStokesWitnessConstructionRouteObstruction.lean}, the definition
\texttt{ExactPairWitnessConstructionRoute} bundles the exact finite-time
witness, the same witness plus a uniform vorticity bound, the same witness plus
BKM envelope data, and the prescribed-velocity exact global-output surface for
one fixed velocity/pressure pair.  The theorem
\texttt{ExactPairWitnessConstructionRoute\_iff\_exists\_finiteTimeWitness}
now makes that bundle exact: for a fixed velocity/pressure pair, the
uniform-vorticity, BKM, and exact-global-output wrappers add no new inhabitants
beyond the underlying exact slab witness.  The corollary
\texttt{not\_ExactPairWitnessConstructionRoute\_of\_not\_exists\_finiteTimeWitness}
then shows that any slabwise exact-witness obstruction already kills that whole
exact-pair repair family.  The boxed classification
\texttt{boxedPartialPeriodizationSteadySeed\_timeOnlyPressure\_ExactPairWitnessConstructionRoute\_iff\_stationaryMomentum\_zeroPressure}
reduces the full fixed-pair route for the manuscript-style boxed steady seed
with a smooth time-only pressure gauge to the same zero-pressure stationary
momentum balance as the original exact witness.  The route is now certified to
depend only on the pressure gauge modulo zero spatial gradient:
\texttt{ExactPairWitnessConstructionRoute.addPressureOfZeroSpatialGradient}
and
\texttt{ExactPairWitnessConstructionRoute\_iff\_addPressureOfZeroSpatialGradient}
show that adding any smooth zero-spatial-gradient pressure field preserves the
whole bundled route, not just one branch.  The route is therefore also blocked
uniformly across the whole harmless pressure-gauge class:
\texttt{exists\_boxedPartialPeriodizationSteadySeed\_zeroSpatialGradientPressure\_ExactPairWitnessConstructionRoute\_iff\_stationaryMomentum\_zeroPressure}
shows that even allowing an arbitrary smooth zero-spatial-gradient pressure
field does not weaken the criterion, so time-only gauges and constant gauges are
already subsumed.  More strongly,
\texttt{boxedPartialPeriodizationSteadySeed\_zeroSpatialGradientPressure\_ExactPairWitnessConstructionRoute\_iff\_zeroPressure}
identifies every individual smooth zero-spatial-gradient gauge with the literal
zero-pressure representative on the boxed steady-seed route.  The fixed-gauge
failure theorem
\texttt{not\_boxedPartialPeriodizationSteadySeed\_zeroSpatialGradientPressure\_ExactPairWitnessConstructionRoute\_of\_stationaryMomentum\_failure}
then turns a single failed zero-pressure stationary residual into a direct
blocker for each such pressure field separately, while the existential failure
corollary
\texttt{not\_exists\_boxedPartialPeriodizationSteadySeed\_zeroSpatialGradientPressure\_ExactPairWitnessConstructionRoute\_of\_stationaryMomentum\_failure}
rules out the whole smooth zero-gradient pressure-repair class as soon as the
zero-pressure stationary residual fails at one slab point.  The older fixed
time-only failure corollary
\texttt{not\_boxedPartialPeriodizationSteadySeed\_timeOnlyPressure\_ExactPairWitnessConstructionRoute\_of\_stationaryMomentum\_failure}
proves that once the zero-pressure stationary residual fails at one slab point,
none of those exact-pair repairs survive for the manuscript-style boxed steady
seed.  The companion theorem
\texttt{boxedPartialPeriodizationSteadySeed\_timeOnlyPressure\_exhibits\_nonMomentumFiniteTimeAndGlobalFields\_without\_exactPairWitnessConstructionRoute}
records that the same seed still carries the smoothness, divergence,
initial-match, slab-energy, and global bounded-energy fields, so the blocker is
at the momentum/solution layer rather than in omitted regularity side data.

The finite-dimensional Galerkin ODE layer now has the corresponding equilibrium
audit.  In \texttt{NavierStokesFiniteModeGalerkinToy.lean}, Lean proves
\texttt{finiteModeProjectedNSRHS\_zero}: at the zero coefficient vector the
projected quadratic right-hand side is exactly the forcing vector.  The theorem
\texttt{finiteModeZeroCurve\_solves\_iff\_forcing\_zero} then says that the
zero coefficient curve solves the projected ODE for all times if and only if
every forcing coefficient vanishes, and
\texttt{not\_forall\_finiteModeZeroCurve\_hasDerivAt\_of\_forcing\_ne\_zero}
turns any one nonzero forcing coefficient into a direct blocker.  Thus a
finite-mode repair cannot use the zero coefficient curve as a silent local
solution unless it has first discharged the forcing-equilibrium condition.
The same audit is now available at arbitrary frozen coefficient states:
\texttt{finiteModeConstantCurve\_solves\_iff\_rhs\_eq\_zero} says that a
constant coefficient curve solves the projected ODE for all times exactly when
the projected RHS vanishes at that coefficient vector, and
\texttt{not\_forall\_finiteModeConstantCurve\_hasDerivAt\_of\_rhs\_ne\_zero}
turns any nonzero projected residual into a direct blocker.
\texttt{finiteModeConstantCurve\_solves\_iff\_forall\_rhs\_component\_eq\_zero}
and
\texttt{not\_forall\_finiteModeConstantCurve\_hasDerivAt\_of\_rhs\_component\_ne\_zero}
record the coordinate form: one visibly nonzero projected mode equation is
already enough to rule out the frozen coefficient curve.
The same file now also contains constructive global solution families for
non-frozen finite-mode ODEs.  For pure forcing coefficients, Lean proves
\texttt{finiteModeProjectedNSRHS\_pureForcing} and the affine solution
\texttt{finiteModeAffineForcingCurve\_solves}, i.e. the coefficient path
\(a_0+t f\) solves \(a'=f\) in every finite mode dimension.  For decoupled
diagonal-linear coefficients, Lean proves
\texttt{finiteModeProjectedNSRHS\_diagonalLinear} and
\texttt{finiteModeDiagonalLinearExpCurve\_solves}, i.e. the componentwise
exponential path \(a_i(t)=a_{0,i}\exp(\lambda_i t)\) solves
\(a_i'=\lambda_i a_i\).  This is deliberately outside the cancelled
anti-profile and shear-residual cells: it grows the reusable finite-mode ODE
base with exact positive examples before any PDE reconstruction claim.

The inhomogeneous diagonal-affine case is now covered in the same concrete
layer.  For coefficients satisfying \(a_i'=f_i+\lambda_i a_i\), Lean proves
\texttt{finiteModeProjectedNSRHS\_diagonalAffine} and
\texttt{finiteModeDiagonalAffineEquilibriumCurve\_solves}: if an explicitly
supplied center \(a^{\mathrm{eq}}\) satisfies
\(f_i+\lambda_i a^{\mathrm{eq}}_i=0\), then
\[
  a_i(t)=a^{\mathrm{eq}}_i+
    (a_{0,i}-a^{\mathrm{eq}}_i)\exp(\lambda_i t)
\]
is a global finite-mode ODE solution.  The theorem avoids a division-based
equilibrium formula and therefore has no nonzero-rate side condition; it is
still an ODE anchor, not a PDE closure or pressure reconstruction theorem.

The same finite-mode toy layer now has an energy derivative hook.  The
definition \texttt{finiteModeEnergy} specializes the existing finite square-sum
\texttt{gamma} to coefficient vectors, while
\texttt{finiteModeEnergy\_hasDerivAt\_of\_projectedNS} proves the exact
identity
\[
  \frac{d}{dt}\sum_i a_i(t)^2 =
    2\sum_i a_i(t)\,\mathrm{RHS}_i(a(t))
\]
for any coefficient path solving the projected ODE.  In the diagonal-linear
case, \texttt{finiteModeDiagonalLinearEnergyRHS\_nonpos} proves dissipativity
from the reusable hypothesis \(\lambda_i\le 0\), and
\texttt{finiteModeDiagonalLinearExpCurve\_energy\_hasDerivAt} gives the exact
energy derivative along the explicit exponential solution.

At the same time, the positive finite-mode Schwartz lane now reaches the
repaired whole-space theorem surfaces themselves once its closure hypotheses are
supplied.  Lean packages
\texttt{ExplicitConcreteNavierStokesGlobalOutput\_twoModeSchwartzVelocity\_of\_momentumEquation},
\texttt{ExplicitFiniteEnergyAdmissibleNavierStokesRegularityClause\_twoModeSchwartzInitialVelocity\_of\_momentumEquation},
and
\texttt{finiteEnergyConcreteNavierStokesGlobalRegularityClause\_twoModeSchwartzInitialVelocity\_of\_momentumEquation}.
So a bounded two-mode Schwartz ansatz with smooth coefficient curves,
divergence-free profiles, and a proved pointwise momentum identity is no longer
merely a finite-time witness package: it now yields an actual explicit global
output and both repaired whole-space clauses on the corresponding finite-energy
initial datum.

Lean now also pushes that same package onto the manuscript-style divergence-free
Schwartz input lane itself.  In
\texttt{NavierStokesFiniteModeSchwartzDivergenceFreeBridge.lean}, the subtype
\texttt{twoModeSchwartzDivergenceFreeInitialVelocity} packages the same initial
slice as an $S_\sigma(\mathbb{R}^3)$ datum, and theorems
\texttt{ExplicitSchwartzDivergenceFreeNavierStokesRegularityClause\_twoModeSchwartzDivergenceFreeInitialVelocity\_of\_momentumEquation}
and
\texttt{StructuredSchwartzDivergenceFreeNavierStokesRegularityClause\_twoModeSchwartzDivergenceFreeInitialVelocity\_of\_momentumEquation}
show that the same bounded two-mode closure hypotheses now feed directly into
the manuscript-style explicit and structured $S_\sigma(\mathbb{R}^3)$ clauses, not
only the ambient repaired whole-space shells.

Lean now pushes that same divergence-free Schwartz package one layer deeper
into the continuation machinery.  In
\texttt{NavierStokesFiniteModeSchwartzBKMBridge.lean}, theorem
\texttt{twoModeSchwartzDivergenceFreeInitialVelocity\_finiteTimeWitness\_of\_momentumEquation}
packages the same ansatz as an exact
\texttt{ExplicitFiniteTimeRegularityWitness}, and theorem
\texttt{twoModeSchwartzDivergenceFreeInitialVelocity\_exhibits\_BKMContinuationPremise\_of\_momentumEquation\_and\_uniformVorticityBound}
shows that any supplied slabwise uniform vorticity bound upgrades that witness
to an explicit integrable BKM envelope.  So the bounded two-mode
divergence-free Schwartz route now reaches the repaired BKM continuation input
surface itself, not only the repaired whole-space theorem shells.

Lean now closes the missing analytic hypothesis for that same bounded two-mode
lane from first principles.  In
\texttt{NavierStokesFiniteModeVorticity.lean}, the explicit constant
\texttt{schwartzInitialVelocityVorticityBound} is shown to dominate the
pointwise vorticity of a time-independent Schwartz seed via the first Schwartz
seminorm.  The single-profile consequences
\texttt{uniformVorticityBoundUpTo\_timeIndependentVelocity} and
\texttt{timeIndependentVelocity\_exhibits\_BKMEnvelope} now package that estimate
directly as a slabwise uniform vorticity bound and an integrable BKM envelope on
every time horizon.  The witness-construction audit now consumes the uniform
half of this package directly:
\texttt{timeIndependentVelocity\_schwartz\_uniformVorticityData\_iff\_stationaryMomentum}
proves that a fixed smooth pressure gives uniform-vorticity-packaged finite-time
data for an exact time-independent Schwartz seed if and only if the stationary
momentum equation already holds on the slab.  Then
\texttt{uniformVorticityBoundUpTo\_twoModeSchwartzVelocity\_of\_abs\_le} and
\texttt{twoModeSchwartzVelocity\_exhibits\_BKMEnvelope\_of\_abs\_le} push that
control to the full bounded two-mode ansatz under the existing coefficient
amplitude bounds.  Consequently,
\texttt{twoModeSchwartzDivergenceFreeInitialVelocity\_exhibits\_BKMContinuationPremise\_of\_momentumEquation}
removes the extra vorticity-bound hypothesis from the BKM bridge: for this
finite-mode Schwartz route, the continuation premise now follows from the
momentum identity, divergence-free initial slice, and bounded amplitudes alone.

The concrete \(\mathbb{R}^3\) vector-calculus layer now has sharper reusable
vorticity norm facts for the elementary shear families.  In
\texttt{VectorCalculusR3.lean}, Lean proves exact pointwise formulas for
\texttt{norm\_spatialVorticity\_linearShearVelocityField},
\texttt{norm\_spatialVorticity\_heatShearVelocityField},
\texttt{norm\_spatialVorticity\_heatShearTransportVelocityField}, and
\texttt{norm\_spatialVorticity\_heatShearTransportFullDriftVelocityField}.
It also proves the corresponding phase-independent envelope bounds with the
explicit damping factor
\[
  \|\omega(t,x)\| \le |a|\,\exp(-(\nu k^2)t)\,|k|
\]
for the heat-shear, transported heat-shear, and transported full-drift
heat-shear branches.  These are not continuation theorems; they are concrete
curl computations that can be reused by the uniform-vorticity and BKM layers
without redoing the trigonometric estimate.
The same concrete layer now records the global integrability obstruction for
arbitrary smooth pressure gradients.  The definition
\texttt{pressureGradientVelocityField} views a smooth pressure as the velocity
field \(u=\nabla_x p\), and
\texttt{spatialVorticity\_pressureGradientVelocityField} proves this velocity
has zero spatial vorticity everywhere.  Consequently
\texttt{not\_exists\_smoothPressure\_pressureGradientVelocityField\_eq\_of\_spatialVorticity\_ne\_zero}
rules out pure pressure-gradient representations for any velocity with one
nonzero-curl point; the regression
\texttt{linear\_shear\_unit\_not\_smooth\_pressure\_gradient\_regression}
applies it to the unit linear shear.  This is not another normal-form lemma for
the finite-mode pressure cell: it is a reusable \(\nabla \times \nabla p=0\)
fact on the \(\mathbb{R}^3\) vector-calculus surface.
The obstruction now also reaches the concrete momentum equation.  Lean defines
\texttt{momentumPressureResidual} as
\(\nu\Delta u-\partial_tu-(u\cdot\nabla)u\), proves that a smooth-pressure
momentum solution must identify this residual with \(\nabla p\), and therefore
proves
\texttt{spatialVorticity\_momentumPressureResidual\_of\_momentumEquation}.
The packaged no-go
\texttt{not\_exists\_smoothPressure\_momentumEquation\_of\_momentumPressureResidual\_vorticity\_ne\_zero}
says that one nonzero-curl residual point is enough to rule out every smooth
pressure.  The concrete theorem
\texttt{not\_exists\_smoothPressure\_momentumEquation\_undampedUnitHeatShearVelocityField\_unitViscosity}
instantiates this on the undamped unit heat-shear profile at unit viscosity,
whose residual is the negative heat-shear profile and has nonzero curl at the
origin.
The same residual audit is now parameterized over wrong viscosities.  Lean
proves
\texttt{momentumPressureResidual\_heatShearVelocityField}, which states that a
heat-shear field evolved at viscosity \(\nu\) but tested at viscosity \(\mu\)
has pressure residual \((\nu-\mu)k^2u_\nu\).  The origin-curl theorem
\texttt{spatialVorticity\_momentumPressureResidual\_heatShearVelocityField\_origin}
then computes the obstruction as \((\mu-\nu)ak^3e_2\), and
\texttt{not\_exists\_smoothPressure\_momentumEquation\_heatShearVelocityField\_wrongViscosity}
rules out every smooth pressure whenever the viscosity mismatch, amplitude, and
frequency are nonzero.  Thus a nondegenerate heat-shear profile cannot be
retrofitted to a different viscosity by choosing a clever pressure.
The same calculation is now stated at the more natural amplitude level.  Lean
defines
\texttt{amplitudeShearVelocityField} for
\[
  u(t,x)=A(t)\sin(kx_1)e_0
\]
and proves slice-congruence lemmas so that the spatial derivative, divergence,
convection, Laplacian, and vorticity calculations reuse the heat-shear slice
without pretending the whole space-time fields are equal.  If \(A\) has sampled
derivative \(A'(t)\), then
\texttt{momentumPressureResidual\_amplitudeShearVelocityField\_slice} computes
the residual amplitude as \(-(A'(t)+\mu k^2A(t))\).  Hence
\texttt{not\_exists\_smoothPressure\_momentumEquation\_amplitudeShearVelocityField\_of\_heatODE\_mismatch\_at}
rules out every smooth pressure at any nonzero-frequency time where the scalar
heat ODE \(A'(t)+\mu k^2A(t)=0\) fails.  This is the durable bottom-up statement:
pressure repair cannot hide a bad time amplitude.
The transported arbitrary-amplitude version is now formalized as well:
\[
  u(t,x)=A(t)\sin(k(x_1-bt))e_0+be_1 .
\]
Lean proves that the phase derivative and the convection term cancel, so the
pressure residual has the same heat-ODE coefficient
\(-(A'(t)+\mu k^2A(t))\).  The theorem
\texttt{not\_exists\_smoothPressure\_momentumEquation\_amplitudeShearTransportVelocityField\_of\_heatODE\_mismatch\_at}
therefore rules out every smooth pressure repair whenever this scalar ODE fails
at a nonzero frequency.  Constant transport does not create an escape route for
bad amplitude dynamics.
The BKM continuation file now consumes that exact surface directly.  The new
\texttt{heatShearExpVorticityEnvelope} definition records the damped envelope,
\texttt{integrableVorticityEnvelopeOn\_heatShearExpVorticityEnvelope} proves it
is integrable on every nonnegative slab when \(\nu\ge 0\), and the three
\texttt{\_has\_expBKMEnvelope} theorems package pure, transported, and
transported full-drift heat shear as explicit BKM-envelope data with the same
coarse scalar integral bound \(T|ak|\).  This keeps the stronger time-dependent
envelope available downstream while preserving the older constant-envelope API.
For the genuinely damped nonzero-frequency branch, Lean now also records the
closed-form scalar budget
\[
  |a|\,|k|\,\frac{1-\exp(-(\nu k^2)T)}{\nu k^2},
\]
via
\texttt{integral\_heatShearExpVorticityEnvelope\_eq\_exactIntegralBound} and
\texttt{integrableVorticityEnvelopeOn\_heatShearExpVorticityEnvelope\_exact}.
The exact package theorems
\texttt{heatShearVelocityField\_has\_exactExpBKMEnvelope},
\texttt{heatShearTransportVelocityField\_has\_exactExpBKMEnvelope}, and
\texttt{heatShearTransportFullDriftVelocityField\_has\_exactExpBKMEnvelope}
therefore remove another artificial slack source from the heat-shear BKM
interface without changing the continuation problem itself.
The corrected energy/BKM bridge now uses that sharper envelope in its heat-shear
obstruction lane.  Lean proves
\texttt{heatShearVelocityField\_has\_expBKMEnvelope\_but\_not\_energyBKMBridge\_kineticIntegrabilityOnSlab}
and
\texttt{heatShearTransportFullDriftVelocityField\_has\_expBKMEnvelope\_but\_not\_energyBKMBridge\_kineticIntegrabilityOnSlab}:
even with the damped BKM envelope packaged explicitly, the pure and transported
full-drift heat-shear candidates cannot enter the energy/BKM bridge because the
bridge demands slab kinetic-energy integrability and these matched whole-space
initial data already fail finite initial kinetic energy.

The corrected energy/continuation bridge now also has the matching arbitrary
pressure middle theorem for bounded two-profile Schwartz velocity fields.  In
\texttt{NavierStokesEnergyContinuationBridge.lean},
\texttt{boundedKineticEnergyUpTo\_of\_add\_scalar\_smul\_schwartz\_of\_abs\_le\_of\_spatialDivergence\_zero\_of\_pressureIntegrable}
assumes the pointwise momentum equation, pressure-pairing integrability, zero
pressure work, and slice-wise divergence-freeness, and then derives
\texttt{boundedKineticEnergyUpTo}.  The velocity side no longer has to be
supplied abstractly: bounded Schwartz profiles give kinetic integrability,
Schwartz slices give the viscous and convection pairing integrability, and
divergence-freeness gives nonlinear energy cancellation.  This keeps the
remaining pressure obligations explicit while avoiding a premature restriction
to affine-plus-localized Schwartz pressures.
The pressure side now has its own bottom-up Schwartz-slice closure.  The new
energy lemmas
\texttt{pressureEnergyPairing\_schwartzPressureSlice},
\texttt{integrable\_pressureEnergyPairing\_of\_schwartzSlice\_schwartzPressureSlice},
and
\texttt{integral\_pressureEnergyPairing\_of\_schwartzSlice\_schwartzPressureSlice\_of\_spatialDivergence\_zero}
show that any time-indexed Schwartz pressure slice has integrable pressure
pairing and zero pressure work against a divergence-free Schwartz velocity
slice.  Consequently,
\texttt{boundedKineticEnergyUpTo\_of\_add\_scalar\_smul\_schwartz\_of\_abs\_le\_of\_spatialDivergence\_zero\_schwartzPressureSlice}
closes the finite-time bounded-energy bridge for bounded two-profile Schwartz
velocities with arbitrary time-varying Schwartz pressure slices.  This is still
a finite-mode/Schwartz result, not a general pressure regularity theorem, but
it removes another avoidable route-specific pressure ansatz.
The pressure-slice bridge now reaches the actual finite-time witness and BKM
premise surfaces for the bounded two-mode divergence-free Schwartz lane.  The
new constructor
\texttt{ExplicitFiniteTimeRegularityWitness.of\_add\_scalar\_smul\_schwartz\_of\_abs\_le\_of\_spatialDivergence\_zero\_schwartzPressureSlice}
fills the \texttt{bounded\_energy\_on} field from the pressure-slice energy
theorem while leaving smoothness, initial matching, incompressibility, and the
pointwise momentum equation explicit.  In
\texttt{NavierStokesFiniteModeSchwartzBKMBridge.lean},
\texttt{twoModeSchwartzDivergenceFreeInitialVelocity\_finiteTimeWitness\_of\_momentumEquation\_schwartzPressureSlice}
then packages the two-mode divergence-free Schwartz ansatz as a finite-time
witness for any smooth time-indexed Schwartz pressure slice satisfying the
momentum equation.  Finally,
\texttt{twoModeSchwartzDivergenceFreeInitialVelocity\_exhibits\_BKMContinuationPremise\_of\_momentumEquation\_schwartzPressureSlice}
combines that witness with the finite-mode vorticity estimate, so bounded
amplitudes supply the integrable BKM envelope internally.  This is still a
premise theorem, not a proof of the BKM continuation clause, but the bounded
two-mode lane no longer needs the affine-plus-localized pressure ansatz to
enter the BKM input surface.
The hidden momentum-equation slot in this pressure-slice route is now exposed
as an explicit residual closure.  In
\texttt{NavierStokesFiniteModeBoundedEnergy.lean},
\texttt{momentumEquation\_twoModeSchwartzVelocity\_schwartzPressureSlice\_of\_explicitClosure}
expands the momentum equation into the two scalar time derivatives, the four
self/mixed convection terms, the pressure-slice gradient, and the two
Laplacian terms.  The BKM bridge theorem
\texttt{twoModeSchwartzDivergenceFreeInitialVelocity\_exhibits\_BKMContinuationPremise\_of\_explicitClosure\_schwartzPressureSlice}
then consumes exactly that closure equation.  Thus the next nontrivial finite
mode burden is not an opaque \texttt{heq} hypothesis but a concrete residual
identity for the chosen profiles, amplitudes, viscosity, and pressure slice.
The pressure-slice route also now has the matching anti-profile cancellation
sanity check.  The theorem
\texttt{explicitClosure\_oneOneAntiProfileSchwartzVelocity\_schwartzPressureSlice\_zeroPressure}
shows that \(f\) paired with \(-f\), with constant unit amplitudes and the
zero Schwartz pressure slice, satisfies the expanded pressure-slice residual
closure for every viscosity.  The BKM bridge theorem
\texttt{twoModeSchwartzDivergenceFreeInitialVelocity\_exhibits\_BKMContinuationPremise\_oneOneAntiProfile\_schwartzPressureSlice\_zeroPressure}
then packages the same cancellation field as BKM-premise data for the zero
initial datum whenever \(f\) is divergence-free and \(\nu\ge 0\).  This is
useful as a route regression because the pressure-slice closure surface now
contains the same zero-flow baseline as the older zero-pressure lane.  It is
not a nontrivial solution: the total velocity and initial datum cancel, so the
non-cancelling finite-mode residual problem remains untouched.
The unit-amplitude restriction has also been removed from that baseline.  The
new equal-amplitude theorem
\texttt{explicitClosure\_equalAmplitudeAntiProfileSchwartzVelocity\_schwartzPressureSlice\_zeroPressure}
allows any smooth scalar amplitude \(a(t)\) in the pair \(a(t)f,
a(t)(-f)\), and the corresponding BKM bridge
\texttt{twoModeSchwartzDivergenceFreeInitialVelocity\_exhibits\_BKMContinuationPremise\_equalAmplitudeAntiProfile\_schwartzPressureSlice\_zeroPressure}
uses a uniform bound on \(a\) to enter the same finite-mode BKM-premise
surface.  This makes the cancellation boundary less brittle: any repair that
tries to exploit time-dependent equal coefficients is already absorbed by the
zero-flow analysis.  It still leaves unequal or non-cancelling profiles as the
real closure burden.
The zero-flow side of that boundary is now exact for nonzero profiles.  Lean
proves
\texttt{antiProfileSchwartzVelocity\_eq\_zero\_iff\_equalAmplitude}: if
\(\exists x,\ f(x)\ne 0\), then the anti-profile velocity
\(a(t)f+b(t)(-f)\) is identically zero if and only if \(a(t)=b(t)\) for every
time.  The initial-slice companion
\texttt{antiProfileSchwartzInitialVelocity\_eq\_zero\_iff\_equalAmplitude}
proves the same statement for \(a_0f+b_0(-f)\), and
\texttt{exists\_antiProfileSchwartzVelocity\_ne\_zero\_of\_exists\_amplitude\_ne}
gives the positive contrapositive: unequal amplitudes somewhere force a
nonzero velocity somewhere.  So the anti-profile escape hatch is now sharply
classified: cancellation is exactly equal-amplitude cancellation.
The BKM bridge now records the same audit boundary at witness level via
\texttt{twoModeSchwartzDivergenceFreeInitialVelocity\_exhibits\_BKMContinuationPremise\_equalAmplitudeAntiProfile\_schwartzPressureSlice\_zeroPressure\_zeroWitness}:
the equal-amplitude anti-profile bridge can be witnessed by zero velocity and
the zero time-indexed Schwartz pressure slice.  Thus the BKM-premise inhabitant
is not merely associated with zero initial data; its actual finite-time
velocity is explicitly the zero flow.  The companion
\texttt{...\_zeroWitness\_zeroEnvelope} theorem fixes the BKM side to the zero
vorticity envelope with zero integral budget, so this cancellation branch now
has no remaining hidden finite-mode BKM content.  The latest
\texttt{...\_canonicalZeroWitness} refinement goes one step further: the
finite-mode witness is propositionally the canonical
\texttt{zeroFiniteTimeRegularityWitness}.  Thus the equal-amplitude
anti-profile branch is not a separate two-mode construction at all.
The same audit now covers the adjacent time-only pressure-gauge repair.  Lean
proves
\texttt{momentumEquation\_equalAmplitudeAntiProfileSchwartzVelocity\_affineTimeOnlyPressure}
for any time-only affine-plus-Schwartz pressure gauge, and the BKM bridge
\texttt{twoModeSchwartzDivergenceFreeInitialVelocity\_exhibits\_BKMContinuationPremise\_equalAmplitudeAntiProfile\_affineTimeOnlyPressure\_zeroWitness\_zeroEnvelope}
packages that branch with zero velocity, the stated gauge pressure, the zero
vorticity envelope, and zero integral budget.  So time-only pressure freedom
does not create a new finite-mode branch; it only changes the pressure gauge on
the zero flow.  The canonical-gauge refinement
\texttt{...\_affineTimeOnlyPressure\_canonicalGaugeWitness} now identifies that
witness with \texttt{zeroFiniteTimeRegularityWitness} after the same
\texttt{addTimeOnlyPressure} operation, making the gauge-only nature of the
repair explicit at the witness level.
The adjacent pure spatial-affine pressure repair is now blocked outright:
\texttt{spatialPressureGradient\_affineAddScalarSchwartzPressure\_spatialAffine}
computes the gradient as the affine coefficient \(c(t)\), and
\texttt{ExplicitFiniteTimeRegularityWitness.zeroVelocity\_spatialAffinePressure\_implies\_zeroAffineCoeffOn}
forces that coefficient to vanish on any certified zero-velocity finite-time
witness.  A nonzero spatial-affine pressure is therefore not another gauge of
the equal-amplitude anti-profile zero branch.  The broader BKM classification
\texttt{BKMData\_zeroVelocity\_implies\_initialVelocity\_zero} also forces the
initial datum to be zero for any zero-velocity BKM-data witness, and
\texttt{not\_exists\_BKMData\_zeroVelocity\_of\_initialVelocity\_ne\_zero}
turns nonzero initial data into a no-go theorem for this branch.  The pressure
classification
\texttt{BKMData\_zeroVelocity\_implies\_spatialPressureGradient\_zero}
removes the ansatz restriction at the packaging layer: any zero-velocity
BKM-data witness forces the pressure's spatial gradient to vanish on the
certified slab.  Its contrapositive
\texttt{not\_exists\_BKMData\_zeroVelocity\_of\_exists\_spatialPressureGradient\_ne\_zero}
turns one nonzero gradient value into a no-go theorem.  The combined theorem
\texttt{BKMData\_zeroVelocity\_implies\_initialVelocity\_zero\_and\_spatialPressureGradient\_zero}
records that both necessary conditions hold simultaneously, and
\texttt{not\_exists\_BKMData\_zeroVelocity\_of\_initialVelocity\_ne\_zero\_or\_exists\_spatialPressureGradient\_ne\_zero}
rules out the zero-flow BKM package whenever either condition fails.  The
\texttt{BKMData\_zeroVelocity\_iff\_initialVelocity\_zero\_and\_spatialPressureGradient\_zero}
theorem makes this boundary exact for smooth pressure fields: zero-flow BKM
data exists precisely for zero initial data and zero spatial pressure gradient
on the certified slab.  The positive direction uses the lower-level zero-flow
finite-time witness classification and adds the zero vorticity envelope, so the
BKM packaging layer is conservative here rather than a repair mechanism.  The
time-only pressure gauge is now exposed as the exact surviving specialization:
\texttt{BKMData\_zeroVelocity\_timeOnlyPressure\_iff\_initialVelocity\_zero}
classifies it by the single condition \(u_0=0\), while
\texttt{BKMData\_zeroVelocity\_timeOnlyPressure} and
\texttt{not\_exists\_BKMData\_zeroVelocity\_timeOnlyPressure\_of\_initialVelocity\_ne\_zero}
give the positive and nonzero-data sides separately.  The spatial-affine
pressure subclass is also exact now:
\texttt{BKMData\_zeroVelocity\_spatialAffinePressure\_iff\_initialVelocity\_zero\_and\_affineCoeff\_zeroOn}
reduces the zero-flow BKM condition to \(u_0=0\) and slabwise \(c(t)=0\).
The constructive theorem
\texttt{BKMData\_zeroVelocity\_spatialAffinePressure\_of\_affineCoeff\_zeroOn}
and the combined no-go
\texttt{not\_exists\_BKMData\_zeroVelocity\_spatialAffinePressure\_of\_initialVelocity\_ne\_zero\_or\_exists\_nonzeroAffineCoeffOn}
keep the concrete affine-plus-Schwartz pressure repair blocked as a named
special case, while still exposing the harmless zero-coefficient side.
The full affine-plus-localized pressure class is now classified too:
\texttt{spatialPressureGradient\_affineAddScalarSchwartzPressure} computes the
forcing term as \(c(t)+\rho(t)\nabla q(x)\), and
\texttt{BKMData\_zeroVelocity\_affineAddScalarSchwartzPressure\_iff\_initialVelocity\_zero\_and\_pressureGradientBalanceOn}
proves that zero-flow BKM data exists exactly when \(u_0=0\) and this balance
vanishes on the certified slab.  The positive balanced case is packaged as
\texttt{BKMData\_zeroVelocity\_affineAddScalarSchwartzPressure\_of\_pressureGradientBalanceOn}.
The obstruction
\texttt{not\_exists\_BKMData\_zeroVelocity\_affineAddScalarSchwartzPressure\_of\_initialVelocity\_ne\_zero\_or\_exists\_nonconstant\_schwartzPressureGradientOn}
records the important repair failure: if \(\rho(t)\neq0\) and the Schwartz
pressure profile has two different spatial-gradient values at that time, no
single affine coefficient \(c(t)\) can cancel the localized pressure force
inside a zero-flow BKM package.  The companion theorem
\texttt{BKMData\_zeroVelocity\_affineAddScalarSchwartzPressure\_implies\_schwartzPressureGradient\_constantOn\_of\_nonzeroAmplitude}
states the exact residual requirement: any such zero-flow BKM package at a
nonzero-amplitude time forces the Schwartz pressure profile to have the same
spatial gradient at every pair of points.  Lean now also proves
\texttt{schwartzMap\_gradient\_constant\_eq\_zero}: a Schwartz pressure profile
on the concrete Navier--Stokes space cannot have a nonzero globally constant
gradient.  The BKM consequences
\texttt{BKMData\_zeroVelocity\_affineAddScalarSchwartzPressure\_implies\_schwartzPressureGradient\_zero\_of\_nonzeroAmplitude}
and
\texttt{BKMData\_zeroVelocity\_affineAddScalarSchwartzPressure\_implies\_affineCoeff\_zero\_of\_nonzeroAmplitude}
therefore force both \(\nabla q=0\) and \(c(t)=0\) at every certified
nonzero-amplitude zero-flow time.
The exact collapsed classification is now also a Lean theorem:
\texttt{BKMData\_zeroVelocity\_affineAddScalarSchwartzPressure\_iff\_initialVelocity\_zero\_and\_affineCoeff\_zeroOn\_and\_schwartzPressureGradient\_zero\_if\_nonzeroAmplitudeOn}.
It says that zero-flow BKM data in the affine-plus-localized Schwartz pressure
class exists precisely when \(u_0=0\), \(c(t)=0\) on the certified slab, and
any nonzero localized amplitude on that slab forces the fixed Schwartz
pressure gradient to vanish everywhere.  This turns the previous one-slice
consequence into the reusable bottom-up interface for later finite-mode
pressure repair checks.
At the raw momentum-equation layer, the equal-amplitude anti-profile branch now
has the matching exact collapse:
\texttt{momentumEquation\_equalAmplitudeAntiProfileSchwartzVelocity\_affineAddScalarSchwartzPressure\_iff\_affineCoeff\_zero\_and\_schwartzPressureGradient\_zero\_if\_nonzeroAmplitude}.
The theorem says that the full pointwise equation for
\texttt{twoModeSchwartzVelocity a a f (-f)} with pressure
\(\langle c(t),x\rangle+\pi(t)+\rho(t)q(x)\) holds precisely when \(c(t)=0\)
for every time, and any nonzero localized amplitude forces the fixed Schwartz
pressure gradient to vanish everywhere.  This keeps the reusable interface
below BKM packaging: the pressure-repair obstruction is already visible in the
momentum equation itself.
The finite-time BKM layer now has the same anti-profile exactness theorem:
\texttt{BKMData\_equalAmplitudeAntiProfileSchwartzVelocity\_affineAddScalarSchwartzPressure\_iff\_affineCoeff\_zeroOn\_and\_schwartzPressureGradient\_zero\_if\_nonzeroAmplitudeOn}.
It classifies BKM data with the cancelled two-mode velocity and full
affine-plus-localized pressure by the slabwise zero-flow pressure conditions,
so later finite-mode repair checks can cite the anti-profile statement directly
without reopening the zero-flow reduction.
The reusable pressure-independent normalization is now available too:
\texttt{BKMData\_equalAmplitudeAntiProfileSchwartzVelocity\_iff\_zeroVelocity}
and
\texttt{BKMData\_equalAmplitudeAntiProfileInitialVelocity\_iff\_zeroVelocity}
identify the cancelled velocity and cancelled initial-datum witness surfaces
with the zero-initial, zero-velocity BKM-data surface for any pressure field.
The same reduction now yields an arbitrary-pressure exact boundary:
\texttt{BKMData\_equalAmplitudeAntiProfileSchwartzVelocity\_iff\_spatialPressureGradient\_zero}
and
\texttt{BKMData\_equalAmplitudeAntiProfileInitialVelocity\_iff\_spatialPressureGradient\_zero}
classify the cancelled anti-profile BKM-data package by zero spatial pressure
gradient on the certified slab for any smooth pressure.  Their direct no-go
forms,
\texttt{not\_exists\_BKMData\_equalAmplitudeAntiProfileSchwartzVelocity\_of\_exists\_spatialPressureGradient\_ne\_zero}
and
\texttt{not\_exists\_BKMData\_equalAmplitudeAntiProfileInitialVelocity\_of\_exists\_spatialPressureGradient\_ne\_zero},
rule out any pressure repair with a nonzero slabwise spatial gradient, without
assuming the affine-plus-localized ansatz.
The arbitrary-pressure no-go now also has a repaired-clause separation form:
\texttt{ExplicitFiniteEnergyBKMContinuationClause\_zero\_and\_not\_exists\_BKMData\_equalAmplitudeAntiProfileSchwartzVelocity\_of\_exists\_spatialPressureGradient\_ne\_zero}
and
\texttt{ExplicitFiniteEnergyBKMContinuationClause\_equalAmplitudeAntiProfileInitialVelocity\_and\_not\_exists\_BKMData\_equalAmplitudeAntiProfileInitialVelocity\_of\_exists\_spatialPressureGradient\_ne\_zero}.
These keep the repaired finite-energy BKM clause available, but deny any
matching anti-profile BKM-data package for pressures with nonzero slabwise
spatial gradient.
The corresponding exact boundary is now also formalized:
\texttt{ExplicitFiniteEnergyBKMContinuationClause\_zero\_and\_BKMData\_equalAmplitudeAntiProfileSchwartzVelocity\_iff\_spatialPressureGradient\_zero}
and
\texttt{ExplicitFiniteEnergyBKMContinuationClause\_equalAmplitudeAntiProfileInitialVelocity\_and\_BKMData\_equalAmplitudeAntiProfileInitialVelocity\_iff\_spatialPressureGradient\_zero}.
They state that the repaired finite-energy BKM clause together with matching
anti-profile BKM data exists exactly when the arbitrary smooth pressure has
zero spatial gradient on the certified slab.  This turns the no-go into a
complete boundary for this finite-mode repair surface.
The raw constructor-input layer has the same arbitrary-pressure boundary:
\texttt{momentumEquation\_equalAmplitudeAntiProfileSchwartzVelocity\_iff\_spatialPressureGradient\_zero}
classifies the full pointwise momentum equation for the cancelled
equal-amplitude anti-profile velocity by zero spatial pressure gradient
everywhere.  The companion theorem
\texttt{ExplicitFiniteEnergyBKMContinuationClause\_equalAmplitudeAntiProfileInitialVelocity\_and\_not\_momentumEquation\_equalAmplitudeAntiProfileSchwartzVelocity\_of\_exists\_spatialPressureGradient\_ne\_zero}
keeps the repaired finite-energy BKM clause for the syntactic anti-profile
initial datum, but rules out the constructor's momentum equation whenever the
arbitrary pressure has a nonzero spatial gradient.  This is a pressure-family
independent obstruction, not another specialization of the concrete
affine-plus-localized ansatz.
Lean now also states this as an exact clause-plus-momentum boundary:
\texttt{ExplicitFiniteEnergyBKMContinuationClause\_equalAmplitudeAntiProfileInitialVelocity\_and\_momentumEquation\_equalAmplitudeAntiProfileSchwartzVelocity\_iff\_spatialPressureGradient\_zero}.
It says that the repaired finite-energy BKM clause over the syntactic
anti-profile initial datum, paired with the raw pointwise momentum equation,
exists exactly when the arbitrary pressure has zero spatial gradient
everywhere.
The combined clause/data/constructor package is now exact too:
\texttt{ExplicitFiniteEnergyBKMContinuationClause\_equalAmplitudeAntiProfileInitialVelocity\_and\_BKMData\_equalAmplitudeAntiProfileInitialVelocity\_and\_momentumEquation\_equalAmplitudeAntiProfileSchwartzVelocity\_iff\_spatialPressureGradient\_zero}.
Thus adding the BKM-data witness does not relax the constructor input: the raw
momentum equation still forces global zero spatial pressure gradient.
The actual finite-time witness surface is now classified directly, without
adding another repaired-clause/BKM-data conjunction:
\texttt{exists\_ExplicitFiniteTimeRegularityWitness\_equalAmplitudeAntiProfileInitialVelocity\_velocity\_pressure\_iff\_spatialPressureGradient\_zero\_on}.
It says that a witness whose velocity is the cancelled anti-profile field and
whose pressure is an arbitrary smooth pressure exists exactly when the pressure
has zero spatial gradient on the certified slab.  The no-go form
\texttt{not\_exists\_ExplicitFiniteTimeRegularityWitness\_equalAmplitudeAntiProfileInitialVelocity\_velocity\_pressure\_of\_exists\_spatialPressureGradient\_ne\_zero\_on}
rules out the witness from one nonzero slabwise pressure-gradient point.
Lean now also proves the corresponding internal-pressure obstruction:
\texttt{ExplicitFiniteTimeRegularityWitness\_equalAmplitudeAntiProfileInitialVelocity\_velocity\_implies\_spatialPressureGradient\_pressure\_zero\_on}.
Thus, once an actual finite-time witness is known to use the cancelled
anti-profile velocity, its own stored pressure field has zero spatial pressure
gradient throughout the certified slab.  The same-witness BKM theorem
\texttt{ExplicitFiniteTimeRegularityWitness\_equalAmplitudeAntiProfileInitialVelocity\_velocity\_has\_zero\_vorticityEnvelope}
then shows that the BKM envelope and integral-budget layer is automatically
inhabited by the zero envelope on that same witness.  The normal-form theorem
\texttt{BKMData\_equalAmplitudeAntiProfileInitialVelocity\_iff\_zeroVorticityEnvelope\_zeroBudget}
strengthens this into an exact BKM-data equivalence: an arbitrary envelope and
budget package exists precisely when the same witness can be presented with the
zero envelope and zero budget.  The no-go theorem
\texttt{not\_exists\_ExplicitFiniteTimeRegularityWitness\_equalAmplitudeAntiProfileInitialVelocity\_velocity\_of\_exists\_pressure\_spatialPressureGradient\_ne\_zero\_on}
is the corresponding negative restatement: such a witness cannot also contain
one nonzero pressure-gradient point on the slab.
The next finite-mode theorem deliberately leaves that cancelled cell:
\texttt{oneOneTwoModeSchwartzVelocity\_nonzero\_inviscidZeroPressure\_BKMPremise\_of\_convectionClosure}
packages a constant-one two-profile branch with \(f+g\) nonzero somewhere,
explicit inviscid convection closure, zero pressure, a finite-time witness,
and an internal BKM envelope.  This gives the grassroots lane a positive
non-cancelling sanity anchor rather than another normal form of the zero
velocity branch.
The same bounded-energy layer now also proves the corresponding
positive-viscosity zero-pressure obstruction.  The theorem
\texttt{explicitClosure\_oneOneTwoModeSchwartzVelocity\_zeroPressure\_of\_momentumEquation}
extracts the condition that the explicit convection residual must equal
\(\nu(\Delta f+\Delta g)\), and
\texttt{not\_momentumEquation\_oneOneTwoMode\_zeroPressure\_of\_inviscidClosure\_lapSum\_ne\_zero}
rules out the zero-pressure positive-viscosity equation whenever \(\nu>0\),
the inviscid closure residual is zero, and the summed profile Laplacian is
nonzero somewhere.  This keeps the non-cancelling anchor honest: moving it
from the inviscid boundary to the Navier--Stokes regime requires an actual
viscous/pressure repair, not another BKM packaging lemma.
That repair is now stated exactly.  Lean proves
\texttt{momentumEquation\_oneOneTwoModeSchwartzVelocity\_iff\_pressureGradient\_lapSum\_of\_inviscidClosure},
which says that the pressure-repaired momentum equation is equivalent to the
single balance \(\nabla p=\nu(\Delta f+\Delta g)\), together with the
single-point mismatch obstruction
\texttt{not\_momentumEquation\_oneOneTwoMode\_of\_inviscidClosure\_pressureGradient\_lapSum\_mismatch}.
For affine-plus-localized pressure fields, the theorem
\texttt{oneOneTwoModeSchwartzVelocity\_nonzero\_posViscosity\_BKMPremise\_of\_inviscidClosure\_pressureGradient}
packages the positive direction: if that exact gradient balance is supplied,
then the nonzero constant-one branch yields a positive-viscosity finite-time
witness and an internal BKM envelope.
The same branch now has a concrete affine-plus-localized obstruction that
does not hand the pressure-repair theorem back as an open task.  The theorem
\texttt{not\_momentumEquation\_oneOneTwoMode\_affineAddScalarSchwartzPressure\_of\_lapSum\_pair\_mismatch}
specializes the arbitrary-pressure balance to
\(\langle c(t),x\rangle+\pi(t)+\rho(t)q(x)\) and subtracts two spatial points.
The affine coefficient \(c(t)\) and time gauge \(\pi(t)\) cancel, so any such
repair must satisfy, for every \(t,x,y\),
\[
  \nu\big((\Delta f+\Delta g)(t,x)-(\Delta f+\Delta g)(t,y)\big)
  =
  \rho(t)\big(\nabla q(x)-\nabla q(y)\big).
\]
Thus a single pair mismatch rules out the full momentum equation for the
structured pressure family.  The corollary
\texttt{not\_momentumEquation\_oneOneTwoMode\_affineAddScalarSchwartzPressure\_of\_zeroLocalizedAmplitude\_lapSum\_pair\_ne}
records the inactive-amplitude case: if \(\rho(t)=0\) and \(\nu>0\), the
summed viscous Laplacian residual must be spatially constant at that time.
Any pair where it differs blocks the affine-plus-localized pressure repair.
The obstruction is now quantified over the amplitude choice itself:
\texttt{not\_exists\_momentumEquation\_oneOneTwoMode\_affineAddScalarSchwartzPressure\_of\_no\_scalar\_lapSum\_pair\_fit}
says that if, at one time, no scalar \(r\) can make the fixed pressure-profile
pair-differences match the viscous Laplacian pair-differences, then no choice
of \(c,\pi,\rho\) in the affine-plus-localized family can satisfy the momentum
equation.  The finite two-pair obstruction
\texttt{not\_exists\_momentumEquation\_oneOneTwoMode\_affineAddScalarSchwartzPressure\_of\_same\_pressureGradient\_pair\_lapSum\_pair\_ne}
records a checkable failure pattern: if two spatial pairs have the same
fixed pressure-gradient pair-difference but different summed viscous Laplacian
residual pair-differences, the same scalar amplitude \(\rho(t)\) cannot
repair both.  The concrete corollary
\texttt{not\_exists\_momentumEquation\_oneOneTwoMode\_affineAddScalarSchwartzPressure\_of\_pressureGradient\_pair\_eq\_lapSum\_pair\_ne}
handles a repeated-gradient witness: if \(\nabla q(x)=\nabla q(y)\) but the
summed viscous Laplacian residual separates \(x\) and \(y\) at the same time,
then no scalar localized amplitude can distinguish the two points.  This
blocks the whole structured pressure family, not only a preselected
\(\rho\)-curve.
The obstruction has also been lifted from the global momentum-equation surface
to the actual finite-time BKM-data surface.  The theorem
\texttt{not\_exists\_BKMData\_oneOneTwoMode\_affineAddScalarSchwartzPressure\_of\_same\_pressureGradient\_pair\_lapSum\_pair\_ne\_on}
uses the witness record's slabwise momentum equation at the four bad pair
points, so the contradiction survives after adding an arbitrary vorticity
envelope and integral budget.  The companion composition theorem
\texttt{oneOneTwoModeSchwartzVelocity\_inviscidBKM\_and\_not\_exists\_posViscosity\_BKMData\_affineAddScalarSchwartzPressure\_of\_same\_pressureGradient\_pair\_lapSum\_pair\_ne\_on}
now records the intended separation in one statement: the non-cancelling
constant-one branch carries a nonzero inviscid BKM-premise package, but under
the finite two-pair mismatch no affine-plus-one-Schwartz pressure repair can
turn that same profile pair into positive-viscosity BKM data on the certified
slab.
Lean now pushes the same obstruction onto the exact witness/global surfaces in
\texttt{NavierStokesFiniteModeAffinePressureExactObstruction.lean}.  The theorem
\texttt{not\_exists\_ExplicitFiniteTimeRegularityWitness\_oneOneTwoModeSchwartzVelocity\_affineAddScalarSchwartzPressure\_of\_same\_pressureGradient\_pair\_lapSum\_pair\_ne\_on}
shows that the BKM no-go was not using extra continuation structure in any
essential way: the constant-one two-mode branch already carries a built-in BKM
envelope, so the same two-pair mismatch blocks every exact finite-time witness
in the full affine-plus-localized pressure family.  Then
\texttt{not\_exists\_ExplicitConcreteNavierStokesGlobalOutputWithVelocityPressure\_oneOneTwoModeSchwartzVelocity\_affineAddScalarSchwartzPressure\_of\_same\_pressureGradient\_pair\_lapSum\_pair\_ne\_on}
lifts the obstruction to exact whole-space outputs by restricting any global
pair to the forbidden slab witness.  The summary theorem
\texttt{oneOneTwoModeSchwartzVelocity\_exactSurfaceSeparation\_of\_same\_pressureGradient\_pair\_lapSum\_pair\_ne\_on}
packages the resulting split: the nonzero inviscid exact witness exists, but
under the same finite two-pair mismatch no affine-plus-localized Schwartz
pressure can promote that velocity to a positive-viscosity exact witness or an
exact whole-space output.
Lean now removes the remaining affine-family dependence at the fixed exact-pair
surface in
\texttt{NavierStokesFiniteModeFixedPressureExactObstruction.lean}.  The theorem
\texttt{ExplicitFiniteTimeRegularityWitness.oneOneTwoModeSchwartzVelocity\_pressureGradient\_eq\_lapSum\_on\_of\_inviscidClosure}
shows that any exact finite-time witness for the inviscidly closed constant-one
two-mode branch must realize the full viscous Laplacian residual as its
pressure gradient at each certified slab point.  Consequently,
\texttt{not\_exists\_ExplicitFiniteTimeRegularityWitness\_oneOneTwoModeSchwartzVelocity\_pressure\_of\_inviscidClosure\_pressureGradient\_lapSum\_mismatch\_on}
and
\texttt{not\_ExplicitConcreteNavierStokesGlobalOutputWithVelocityPressure\_oneOneTwoModeSchwartzVelocity\_pressure\_of\_inviscidClosure\_pressureGradient\_lapSum\_mismatch\_on}
block every preselected smooth pressure field with even one mismatch point on
the slab, not only the affine-plus-one-Schwartz subclass.  The concrete
zero-pressure corollaries record the manuscript-relevant instance: if the
summed viscous Laplacian residual is nonzero at one slab point, then the old
inviscid zero-pressure branch cannot be reused as a positive-viscosity exact
finite-time witness or exact whole-space output.
Lean now also records the converse side of that fixed-pressure boundary in
\texttt{NavierStokesFiniteModeFixedPressureExactClassification.lean}.  First,
the generic theorem
\texttt{timeIndependentVelocity\_globalOutputWithVelocityPressure\_iff\_stationaryMomentum}
is added to
\texttt{NavierStokesWitnessConstructionGlobalOutputObstruction.lean}: for any
time-independent seed of finite energy, exact whole-space output with a fixed
pressure is equivalent to the stationary momentum balance.  Specializing that
general fact to the inviscidly closed constant-one branch yields the exact
classification theorems
\texttt{oneOneTwoModeSchwartzVelocity\_finiteTimeWitness\_pressure\_iff\_inviscidClosure\_pressureGradient\_lapSum\_on}
and
\texttt{oneOneTwoModeSchwartzVelocity\_globalOutputWithVelocityPressure\_iff\_inviscidClosure\_pressureGradient\_lapSum}.
They say that for a fixed smooth pressure \(p\), the same one-one velocity is
an exact slab witness exactly when \(\nabla p\) matches the viscous Laplacian
residual on the certified slab, and it is an exact whole-space output exactly
when the same identity holds everywhere.  The manuscript-facing affine-plus-
localized corollaries
\texttt{oneOneTwoModeSchwartzVelocity\_affineAddScalarSchwartzPressure\_finiteTimeWitness\_iff\_inviscidClosure\_pressureGradient\_lapSum\_on}
and
\texttt{oneOneTwoModeSchwartzVelocity\_affineAddScalarSchwartzPressure\_globalOutput\_iff\_inviscidClosure\_pressureGradient\_lapSum}
put the same boundary directly on the author's preferred pressure family.  So
the fixed-pressure layer is now honest in both directions: a mismatch point
blocks the route, while an exact compensating pressure-gradient identity is
already enough to construct the exact witness/global pair.
Lean now also records the converse side of that BKM boundary in
\texttt{NavierStokesFiniteModeFixedPressureBKMClassification.lean}.  The
generic theorem
\texttt{oneOneTwoModeSchwartzVelocity\_BKMData\_pressure\_iff\_inviscidClosure\_pressureGradient\_lapSum\_on}
says that for a fixed smooth pressure \(p\), the inviscidly closed constant-one
two-mode branch admits a BKM-packaged witness with that exact same
velocity/pressure pair if and only if \(\nabla p\) matches the viscous
Laplacian residual on the certified slab.  So the BKM envelope and integral
budget add no new fixed-pressure escape hatch beyond the already formalized
exact witness surface.  The manuscript-facing affine specialization
\texttt{oneOneTwoModeSchwartzVelocity\_affineAddScalarSchwartzPressure\_BKMData\_iff\_inviscidClosure\_pressureGradient\_lapSum\_on}
puts the same exact boundary on the author's preferred affine-plus-localized
pressure family, and the positive-viscosity zero-pressure corollary
\texttt{oneOneTwoModeSchwartzVelocity\_zeroPressure\_BKMData\_iff\_inviscidClosure\_lapSum\_zero\_on}
shows that the old inviscid zero-pressure branch survives on the BKM surface
exactly in the degenerate case where the summed viscous Laplacian residual
already vanishes throughout the slab.
The same fixed-pressure obstruction also survives the BKM packaging layer in
\texttt{NavierStokesFiniteModeFixedPressureBKMObstruction.lean}.  The theorem
\texttt{not\_exists\_BKMData\_oneOneTwoModeSchwartzVelocity\_pressure\_of\_inviscidClosure\_pressureGradient\_lapSum\_mismatch\_on}
simply composes the exact fixed-pressure witness no-go with the generic lift
\texttt{not\_exists\_BKMData\_of\_not\_exists\_finiteTimeWitness}: once the
constant-one two-mode branch is inviscidly closed, adding a vorticity envelope
and integral budget does not relax the pointwise pressure-gradient matching
obligation.  The zero-pressure corollary
\texttt{not\_exists\_BKMData\_oneOneTwoModeSchwartzVelocity\_zeroPressure\_of\_inviscidClosure\_lapSum\_ne\_zero\_on}
therefore rules out the manuscript's old inviscid zero-pressure branch at
positive viscosity on the repaired BKM surface as soon as one nonzero summed
viscous Laplacian residual point is exhibited.  The summary theorem
\texttt{oneOneTwoModeSchwartzVelocity\_inviscidZeroPressure\_BKMPremise\_and\_not\_exists\_posViscosity\_BKMData\_zeroPressure\_of\_inviscidClosure\_lapSum\_ne\_zero\_on}
pairs that no-go with the already available nonzero inviscid BKM premise.
The same fixed-pressure obstruction also survives the older
uniform-vorticity-packaged surface in
\texttt{NavierStokesFiniteModeFixedPressureUniformClassification.lean}.  The
generic theorem
\texttt{oneOneTwoModeSchwartzVelocity\_uniformVorticityData\_pressure\_iff\_inviscidClosure\_pressureGradient\_lapSum\_on}
says that for a fixed smooth pressure \(p\), the inviscidly closed constant-one
two-mode branch admits a uniform-vorticity-packaged witness with that same
velocity/pressure pair if and only if \(\nabla p\) matches the viscous
Laplacian residual on the certified slab.  So a slabwise vorticity bound adds
no extra fixed-pressure escape hatch beyond the exact witness surface.  The
manuscript-facing affine specialization
\texttt{oneOneTwoModeSchwartzVelocity\_affineAddScalarSchwartzPressure\_uniformVorticityData\_iff\_inviscidClosure\_pressureGradient\_lapSum\_on}
puts the same boundary on the author's preferred affine-plus-localized pressure
family, and the positive-viscosity zero-pressure corollary
\texttt{oneOneTwoModeSchwartzVelocity\_zeroPressure\_uniformVorticityData\_iff\_inviscidClosure\_lapSum\_zero\_on}
shows that the old inviscid zero-pressure branch survives on the
uniform-vorticity surface exactly in the degenerate case where the summed
viscous Laplacian residual already vanishes throughout the slab.
The same fixed-pressure obstruction also survives the older
uniform-vorticity-packaged surface in
\texttt{NavierStokesFiniteModeFixedPressureUniformObstruction.lean}.  The
theorem
\texttt{not\_exists\_uniformVorticityData\_oneOneTwoModeSchwartzVelocity\_pressure\_of\_inviscidClosure\_pressureGradient\_lapSum\_mismatch\_on}
again just composes the exact fixed-pressure witness no-go with the generic
lift
\texttt{not\_exists\_uniformVorticityData\_of\_not\_exists\_finiteTimeWitness}:
an added slab vorticity bound does not relax the pointwise pressure-gradient
matching condition.  Its zero-pressure corollary
\texttt{not\_exists\_uniformVorticityData\_oneOneTwoModeSchwartzVelocity\_zeroPressure\_of\_inviscidClosure\_lapSum\_ne\_zero\_on}
therefore blocks the old inviscid zero-pressure branch at positive viscosity
on the uniform-vorticity surface too, and the summary theorem
\texttt{oneOneTwoModeSchwartzVelocity\_inviscidZeroPressure\_uniformVorticityPremise\_and\_not\_exists\_posViscosity\_uniformVorticityData\_zeroPressure\_of\_inviscidClosure\_lapSum\_ne\_zero\_on}
pairs that no-go with the explicit inviscid witness and its slabwise
vorticity bound.
Lean now also records the converse pressure-gauge rigidity in
\texttt{NavierStokesPressureGaugeRigidity.lean}.  The slab-local transport
lemma
\texttt{ExplicitFiniteTimeRegularityWitness.addPressureOfZeroSpatialGradientOn}
shows that an exact finite-time witness is unchanged after adding a pressure
correction whose spatial gradient vanishes only on the certified slab, not
necessarily globally.  More importantly,
\texttt{exists\_finiteTimeWitness\_with\_sameVelocity\_pressure\_iff\_pressureDifference\_zeroGradient\_on}
says that once an exact witness already exists for a fixed velocity \(u\) and
pressure \(p\), another pressure \(q\) yields the same velocity on the same
slab if and only if \(\nabla (q-p)\) vanishes on that slab.  The whole-space
analogue
\texttt{ExplicitConcreteNavierStokesGlobalOutputWithVelocityPressure\_iff\_pressureDifference\_zeroGradient}
gives the same rigidity for exact prescribed-velocity global outputs.  So on
the one-one finite-mode branch, the repaired fixed-pressure realizations
isolated above are already unique modulo time-only gauge: changing the
pressure presentation away from its slabwise spatial gradient does not create
a genuinely new exact repair.
The follow-up file
\texttt{NavierStokesSpatialAffinePressureGaugeRigidity.lean} lowers that
generic statement to the manuscript's spatial-affine pressure lane
\(\texttt{affineAddScalarSchwartzPressure}\; c\; \pi\; 0\; q\).  The theorem
\texttt{exists\_finiteTimeWitness\_with\_sameVelocity\_spatialAffinePressure\_iff\_affineCoeff\_eq\_on}
shows that on a fixed exact slab surface, two such pressures produce the same
velocity if and only if their spatial affine coefficients \(c\) agree on the
certified slab; the time-only scalar term \(\pi\) and idle Schwartz profile
\(q\) are pure gauge.  The whole-space analogue
\texttt{ExplicitConcreteNavierStokesGlobalOutputWithVelocityPressure\_spatialAffinePressure\_iff\_affineCoeff\_eq}
gives the same collapse for exact prescribed-velocity global outputs.  Finally,
the concrete theorem
\texttt{oneOneTwoModeSchwartzVelocity\_spatialAffinePressure\_repair\_iff\_affineCoeff\_eq\_on}
shows that on the inviscidly closed constant-one two-mode branch, once one
spatial-affine repair coefficient closes the slab, every other zero-localized
affine-plus-Schwartz pressure repairs the same velocity exactly when it has the
same coefficient \(c\).  So the manuscript's apparent \((\pi,q)\) repair family
does not add genuinely new exact solutions.
The stronger pressure-agnostic version is now available in
\texttt{NavierStokesFiniteModeOneOneResidualCurlObstruction.lean}.  There
\texttt{momentumPressureResidual\_oneOneTwoModeSchwartzVelocity\_of\_inviscidClosure}
identifies the positive-viscosity exact pressure residual of the inviscidly
closed constant-one branch with the pure viscous Laplacian residual field.
Consequently, if that residual field has nonzero vorticity at one slab point,
then no pressure choice can rescue the same velocity on any exact surface:
\texttt{not\_exists\_ExplicitFiniteTimeRegularityWitness\_oneOneTwoModeSchwartzVelocity\_of\_inviscidClosure\_residualVorticity\_ne\_zero\_on},
\texttt{not\_exists\_BKMData\_oneOneTwoModeSchwartzVelocity\_of\_inviscidClosure\_residualVorticity\_ne\_zero\_on},
\texttt{not\_exists\_uniformVorticityData\_oneOneTwoModeSchwartzVelocity\_of\_inviscidClosure\_residualVorticity\_ne\_zero\_on},
and
\texttt{not\_ExplicitConcreteNavierStokesGlobalOutputWithVelocity\_oneOneTwoModeSchwartzVelocity\_of\_inviscidClosure\_residualVorticity\_ne\_zero\_on}
rule out exact finite-time witnesses, BKM data, uniform-vorticity data, and
prescribed-velocity whole-space outputs alike.  The summary theorem
\texttt{oneOneTwoModeSchwartzVelocity\_pressureAgnosticSurfaceSeparation\_of\_inviscidClosure\_residualVorticity\_ne\_zero\_on}
packages this as a genuine pressure-repair blocker: the same nonzero inviscid
zero-pressure witness still carries explicit BKM and uniform-vorticity data,
but once one residual-curl witness point is exhibited, every positive-viscosity
pressure repair strategy for that velocity is blocked at once.
Complementing that velocity-specific obstruction, the new file
\texttt{NavierStokesSchwartzDirectionalRigidity.lean} proves a structural
Schwartz-lane no-go for one-dimensional geometry itself.  The predicate
\texttt{TranslationInvariantAlong} packages invariance along a fixed nonzero
direction in \(\mathbb{R}^3\),
\texttt{translationInvariantAlong\_of\_fderiv\_eq\_zero\_along} derives that
invariance from vanishing directional Fr\'echet derivative, and
\texttt{schwartzMap\_eq\_zero\_of\_translationInvariantAlong} together with
\texttt{schwartzMap\_eq\_zero\_of\_fderiv\_eq\_zero\_along} show that any such
Schwartz profile must vanish identically.  The concrete initial-data
specializations
\texttt{schwartzInitialVelocity\_eq\_zero\_of\_translationInvariantAlong} and
\texttt{schwartzInitialVelocity\_eq\_zero\_of\_fderiv\_eq\_zero\_along} make
the route consequence explicit: the nondecaying line-invariant shear templates
used elsewhere in the manuscript cannot be transferred unchanged into the
finite-energy Schwartz lane, because any nonzero Schwartz candidate must break
that line invariance before the Navier--Stokes equation is even checked.
The follow-up file
\texttt{NavierStokesSchwartzHeatShearObstruction.lean} applies this structural
rigidity to the explicit heat-shear datum
\(x \mapsto (a \sin(k x_1),0,0)\).  The theorem
\texttt{translationInvariantAlong\_heatShearInitialVelocity\_streamwise} records
that this datum is invariant under translation in the streamwise \(x_0\)
direction, \texttt{heatShearInitialVelocity\_ne\_zero\_of\_amplitude\_ne\_zero\_of\_wavenumber\_ne\_zero}
shows by evaluation at \(x_1 = (\pi/2)/k\) that the datum is genuinely nonzero
when \(a \neq 0\) and \(k \neq 0\), and
\texttt{not\_exists\_schwartzInitialVelocity\_eq\_heatShearInitialVelocity\_of\_amplitude\_ne\_zero\_of\_wavenumber\_ne\_zero}
combines the two facts to rule out any Schwartz initial-velocity realization of
that explicit heat-shear profile.  So this route is blocked before any
pressure, witness, or continuation data are introduced: the plain nondecaying
heat-shear geometry itself does not fit inside the Schwartz input lane unless
it collapses to the zero field.
The companion file
\texttt{NavierStokesSchwartzTransportedHeatShearObstruction.lean} pushes the
same argument through the broader transported full-drift family
\(x \mapsto (d + a \sin(k x_1), b, c)\).  The theorem
\texttt{translationInvariantAlong\_heatShearTransportFullDriftInitialVelocity\_streamwise}
shows that these data remain streamwise-translation-invariant, while
\texttt{heatShearTransportFullDriftInitialVelocity\_eq\_zero\_iff} proves the
exact degeneration criterion \(b = d = c = 0\) and \(a k = 0\).  From this,
\texttt{not\_exists\_schwartzInitialVelocity\_eq\_heatShearTransportFullDriftInitialVelocity\_of\_nonzero\_transport\_or\_streamwiseDrift\_or\_verticalDrift\_or\_amplitude\_mul\_wavenumber\_ne\_zero}
rules out every nondegenerate transported full-drift profile from the Schwartz
lane, and the transported-heat-shear corollary
\texttt{not\_exists\_schwartzInitialVelocity\_eq\_heatShearTransportInitialVelocity\_of\_transport\_ne\_zero\_or\_amplitude\_mul\_wavenumber\_ne\_zero}
does the same for the manuscript's transported branch.  So the Schwartz-lane
obstruction is no longer limited to the base oscillatory slice: adding the
constant transport, streamwise drift, or vertical drift components still does
not rescue these one-dimensional heat-shear templates.
The same Schwartz-lane rigidity now blocks the affine linear-shear family
\(x \mapsto (a x_1,b,c)\).  In
\texttt{NavierStokesSchwartzLinearShearObstruction.lean},
\texttt{translationInvariantAlong\_linearShearFullDriftInitialVelocity\_streamwise}
shows that the full-drift affine datum is again streamwise-translation
invariant, while
\texttt{linearShearFullDriftInitialVelocity\_eq\_zero\_iff} proves the exact
degeneration criterion \(a=b=c=0\).  The theorem
\texttt{not\_exists\_schwartzInitialVelocity\_eq\_linearShearFullDriftInitialVelocity\_of\_parameter\_ne\_zero}
then excludes every nontrivial affine full-drift linear-shear datum from the
Schwartz lane, and the base, horizontal-drift, and vertical-drift corollaries
do the same for the standard subfamilies.  So the nondecaying one-dimensional
obstruction is no longer tied to oscillatory heat-shear templates: even the
affine shear branch remains outside the Schwartz input surface unless all
parameters collapse to zero.
The same line-invariance idea now has a stronger whole-space finite-energy
formulation.  In
\texttt{NavierStokesStreamwiseInvariantFiniteEnergyObstruction.lean}, the
determinant-\(2\) anisotropic map
\texttt{streamwiseScaling} doubles only the streamwise \(x_0\) direction, while
\texttt{not\_integrable\_of\_continuous\_nonneg\_translationInvariantAlong\_streamwise\_of\_exists\_ne\_zero\_point}
shows that any nonnegative continuous density invariant under streamwise
translation cannot be integrable unless it vanishes identically.  Applied to
initial kinetic energy, this yields
\texttt{not\_finiteInitialKineticEnergy\_of\_translationInvariantAlong\_streamwise\_of\_ne\_zero}
and
\texttt{not\_nonempty\_FiniteEnergyAdmissibleNavierStokesProblemData\_of\_translationInvariantAlong\_streamwise\_of\_ne\_zero}:
every nonzero smooth streamwise-invariant initial velocity already lies outside
the repaired finite-energy whole-space input domain.  So for the current
manuscript-facing one-dimensional shear families, the obstruction is not merely
that Schwartz decay fails; the finite-energy hypothesis itself already rules
out the unchanged streamwise-invariant geometry.
The follow-up file
\texttt{NavierStokesHeatShearFiniteEnergyObstruction.lean} now applies that
generic whole-space argument directly to the manuscript heat-shear families.
The theorem
\texttt{not\_finiteInitialKineticEnergy\_heatShearFullDriftInitialVelocity\_of\_streamwiseDrift\_ne\_zero\_or\_verticalDrift\_ne\_zero\_or\_amplitude\_mul\_wavenumber\_ne\_zero}
shows that every nondegenerate full-drift profile
\(x \mapsto (d + a \sin(kx_1),0,c)\) already has infinite initial kinetic
energy, while
\texttt{not\_finiteInitialKineticEnergy\_heatShearTransportFullDriftInitialVelocity\_of\_transport\_ne\_zero\_or\_streamwiseDrift\_ne\_zero\_or\_verticalDrift\_ne\_zero\_or\_amplitude\_mul\_wavenumber\_ne\_zero}
extends the same exclusion to the transported family
\(x \mapsto (d + a \sin(kx_1),b,c)\).  Their
\texttt{not\_nonempty\_FiniteEnergyAdmissibleNavierStokesProblemData\_\dots}
corollaries sharpen the route statement: once any of the manuscript's heat
shear parameters is genuinely switched on, the repaired finite-energy input
type itself is empty on that seed family, before any pressure repair,
continuation package, or exact witness is considered.
The same whole-space finite-energy strategy now also covers the affine
linear-shear branch.  In
\texttt{NavierStokesLinearShearFiniteEnergyObstruction.lean}, the theorem
\texttt{not\_finiteInitialKineticEnergy\_linearShearFullDriftInitialVelocity\_of\_parameter\_ne\_zero}
shows that every nontrivial affine full-drift profile
\(x \mapsto (a x_1,b,c)\) already has infinite initial kinetic energy, and
\texttt{not\_nonempty\_FiniteEnergyAdmissibleNavierStokesProblemData\_linearShearFullDriftInitialVelocity\_of\_parameter\_ne\_zero}
turns that directly into emptiness of the repaired finite-energy admissible
input type on the same family.  The base, vertical-drift, and
horizontal-drift corollaries show that adding constant drift components does
not rescue affine shear either.  So the streamwise-invariant finite-energy
blocker now covers both canonical one-dimensional nondecaying seed styles used
in this route: oscillatory heat shear and affine linear shear.
The same idea has now been factored into a reusable finite-time witness theorem
that no longer mentions any special pressure family.  In
\texttt{NavierStokesUniformVorticityContinuationTarget.lean},
\texttt{not\_exists\_ExplicitFiniteTimeRegularityWitness\_velocity\_of\_momentumPressureResidual\_vorticity\_ne\_zero\_on}
says that a velocity whose exact pressure residual
\(\nu\Delta u-\partial_tu-(u\cdot\nabla)u\) has nonzero curl at one point of the
slab cannot be the velocity field of any smooth finite-time witness on that
slab.  The proof is a direct composition of the witness momentum equation with
the already proved curl-of-gradient identity.  The BKM-level corollary
\texttt{not\_exists\_BKMData\_velocity\_of\_momentumPressureResidual\_vorticity\_ne\_zero\_on}
then shows that adding a vorticity envelope and integral budget cannot rescue
such a velocity.  This is grassroots infrastructure rather than a new analytic
continuation theorem: it exposes a minimal necessary condition every future
finite-time/BKM witness construction must satisfy.
The same necessary condition is now available as positive interface lemmas,
not only as a no-existence theorem:
\texttt{ExplicitFiniteTimeRegularityWitness.spatialVorticity\_momentumPressureResidual\_zero\_on}
and
\texttt{BKMData\_momentumPressureResidual\_vorticity\_zero\_on}
state directly that certified finite-time and BKM packages have curl-free exact
pressure residuals on the certified slab.  Their velocity-equality variants
let later concrete ansatz calculations plug in by identifying the candidate
velocity and producing a single nonzero residual-curl point.
This interface now also reaches the global-output layer.  The wrapper
\texttt{ExplicitConcreteNavierStokesGlobalOutputWithVelocity} fixes the
candidate velocity while remaining propositionally equivalent to the old
global-output definition after existentially quantifying that velocity.  Its
theorem
\texttt{ExplicitConcreteNavierStokesGlobalOutputWithVelocity.spatialVorticity\_momentumPressureResidual\_zero}
states that any such global output has curl-free exact pressure residual
everywhere, and
\texttt{not\_ExplicitConcreteNavierStokesGlobalOutputWithVelocity\_of\_momentumPressureResidual\_vorticity\_ne\_zero}
turns any nonzero residual-curl computation into a direct prescribed-velocity
global-output obstruction.
The transported arbitrary-amplitude shear computation has now been composed
through this whole target stack in
\texttt{NavierStokesTransportedShearResidualCurlObstruction.lean}.  The theorem
\texttt{not\_exists\_ExplicitFiniteTimeRegularityWitness\_velocity\_amplitudeShearTransportVelocityField\_of\_heatODE\_mismatch\_at}
rules out a finite-time witness whenever \(A'(t)+\nu k^2A(t)\neq0\) at a
certified slab time and \(k\neq0\).  The BKM theorem
\texttt{not\_exists\_BKMData\_velocity\_amplitudeShearTransportVelocityField\_of\_heatODE\_mismatch\_at}
shows that the envelope/integral layer adds no repair, and
\texttt{not\_ExplicitConcreteNavierStokesGlobalOutputWithVelocity\_amplitudeShearTransportVelocityField\_of\_heatODE\_mismatch\_at}
rules out the same candidate velocity globally.  Thus the earlier phase-point
curl calculation is no longer only a smooth-pressure no-go; it is a concrete
finite-time, BKM, and global-output obstruction.
The new file
\texttt{NavierStokesWitnessConstructionVelocityRouteObstruction.lean} now
packages that pressure-agnostic witness-construction lane itself.  Its
definition \texttt{VelocityWitnessConstructionRoute} bundles four currently
formalized repair surfaces for one prescribed velocity: exact finite-time
witness, the same witness plus a uniform slab vorticity bound, the same witness
plus BKM envelope data, and prescribed-velocity exact global output.  The
classification theorem
\texttt{VelocityWitnessConstructionRoute\_iff\_exists\_finiteTimeWitness\_velocity}
shows that this whole pressure-agnostic bundle is conservative: the extra
uniform-vorticity, BKM, and global-output wrappers add no new velocity
inhabitants beyond the underlying exact slab witness.  The generic theorem
\texttt{not\_exists\_uniformVorticityData\_velocity\_of\_momentumPressureResidual\_vorticity\_ne\_zero\_on}
fills the remaining older uniform-vorticity branch on the same residual-curl
surface, and
\texttt{not\_VelocityWitnessConstructionRoute\_amplitudeShearTransportVelocityField\_of\_heatODE\_mismatch\_at}
turns the transported-shear phase-point heat-ODE defect into a direct blocker
for the entire pressure-agnostic route at once.  So this concrete residual-curl
example now closes the exact finite-time, uniform-vorticity, BKM, and
prescribed-global witness-construction surfaces together, without choosing any
pressure ansatz.
The same bundled route is now specialized to the constant-one one-one
finite-mode branch in
\texttt{NavierStokesFiniteModeVelocityRouteObstruction.lean}.  The theorem
\texttt{oneOneTwoModeSchwartzVelocity\_nonzero\_and\_inviscidZeroPressure\_exhibits\_velocityRoute\_of\_convectionClosure}
packages the previously constructed inviscid zero-pressure witness as a
genuine inhabitant of the bundled pressure-agnostic route at viscosity \(0\),
so this branch is not merely nonempty on an exact surface; it already occupies
the whole currently formalized witness-construction lane.  Conversely
\texttt{not\_VelocityWitnessConstructionRoute\_oneOneTwoModeSchwartzVelocity\_of\_inviscidClosure\_residualVorticity\_ne\_zero\_on}
shows that one nonzero vorticity point of the viscous Laplacian residual
blocks that same bundled route at viscosity \(\nu\) without choosing any
pressure field.  The combined theorem
\texttt{oneOneTwoModeSchwartzVelocity\_nonzero\_inviscidVelocityRoute\_and\_not\_posViscosity\_velocityRoute\_of\_inviscidClosure\_residualVorticity\_ne\_zero\_on}
therefore gives a second concrete pressure-agnostic route separation beyond
transported shear.
Lean now also records the constructive side of that same route boundary in
\texttt{NavierStokesFiniteModeVelocityRouteConstruction.lean}.  The generic
theorem
\texttt{twoModeSchwartzDivergenceFreeInitialVelocity\_exhibits\_exact\_BKMContinuationPremise\_and\_velocityRoute\_of\_momentumEquation}
states that once a bounded two-mode divergence-free Schwartz ansatz satisfies
the full pointwise Navier--Stokes momentum equation, Lean already has one
explicit finite-time witness with the same velocity and
affine-plus-localized-Schwartz pressure, together with a BKM envelope for that
same witness, and hence an inhabitant of the full pressure-agnostic velocity
route.  The corollary
\texttt{twoModeSchwartzDivergenceFreeInitialVelocity\_exhibits\_velocityRoute\_of\_momentumEquation}
isolates the route inhabitant itself.  More concretely,
\texttt{oneOneTwoModeSchwartzVelocity\_nonzero\_and\_posViscosity\_pressureRepair\_exhibits\_exact\_BKMContinuationPremise\_and\_velocityRoute\_of\_inviscidClosure\_pressureGradient}
shows that the nonzero constant-one branch re-enters the exact/BKM/route stack
at positive viscosity as soon as one supplies the precise compensating pressure
gradient that cancels the viscous Laplacian residual.  So the current
finite-mode route obstruction is paired with an equally concrete constructive
theorem: when the full repair equation is met, the branch really does occupy
the bundled route.
For the original affine-plus-localized Schwartz pressure family, Lean now
states the same paired boundary with the concrete collapsed pressure
conditions:
\texttt{ExplicitFiniteEnergyBKMContinuationClause\_zero\_and\_BKMData\_equalAmplitudeAntiProfileSchwartzVelocity\_affineAddScalarSchwartzPressure\_iff\_affineCoeff\_zeroOn\_and\_schwartzPressureGradient\_zero\_if\_nonzeroAmplitudeOn}
and
\texttt{ExplicitFiniteEnergyBKMContinuationClause\_equalAmplitudeAntiProfileInitialVelocity\_and\_BKMData\_equalAmplitudeAntiProfileInitialVelocity\_affineAddScalarSchwartzPressure\_iff\_affineCoeff\_zeroOn\_and\_schwartzPressureGradient\_zero\_if\_nonzeroAmplitudeOn}.
They say that the repaired clause plus matching BKM data exists exactly when
the affine coefficient vanishes on the certified slab and every active
localized amplitude forces the fixed Schwartz pressure gradient to vanish.
So the concrete pressure ansatz has no separate repair space beyond the
already-collapsed BKM-data classification.
Lean now also states that exact classification over the syntactic
anti-profile initial witness type:
\texttt{BKMData\_equalAmplitudeAntiProfileInitialVelocity\_affineAddScalarSchwartzPressure\_iff\_affineCoeff\_zeroOn\_and\_schwartzPressureGradient\_zero\_if\_nonzeroAmplitudeOn}.
The witness type mentions
\(\texttt{twoModeSchwartzInitialVelocity}\ (a\,0)\ (a\,0)\ f\ (-f)\), but the
classification remains the same slabwise zero-flow pressure condition.  Its
direct no-go form,
\texttt{not\_exists\_BKMData\_equalAmplitudeAntiProfileInitialVelocity\_affineAddScalarSchwartzPressure\_of\_exists\_affineCoeff\_ne\_zero\_on\_or\_active\_schwartzPressureGradient\_ne\_zero},
rules out the matching syntactic witness package under the same bad-pressure
hypotheses.
The same classification is now available in a direct negative form,
\texttt{not\_exists\_BKMData\_equalAmplitudeAntiProfileSchwartzVelocity\_affineAddScalarSchwartzPressure\_of\_exists\_affineCoeff\_ne\_zero\_on\_or\_active\_schwartzPressureGradient\_ne\_zero}.
It rules out the BKM-data package whenever the affine coefficient is nonzero on
the slab, or whenever an active localized pressure amplitude is coupled to a
nonzero fixed Schwartz pressure gradient.
Lean now also records the clause/data separation theorem
\texttt{ExplicitFiniteEnergyBKMContinuationClause\_zero\_and\_not\_exists\_BKMData\_equalAmplitudeAntiProfileSchwartzVelocity\_affineAddScalarSchwartzPressure\_of\_exists\_affineCoeff\_ne\_zero\_on\_or\_active\_schwartzPressureGradient\_ne\_zero}.
It shows that the repaired finite-energy BKM continuation clause at the zero
datum can hold while the structured anti-profile BKM-data package with the bad
affine-plus-localized pressure is impossible.  This keeps later uses honest:
the bare continuation clause is not evidence that the proposed pressure ansatz
has been certified.
The syntactic anti-profile initial-datum version is now recorded as
\texttt{ExplicitFiniteEnergyBKMContinuationClause\_equalAmplitudeAntiProfileInitialVelocity\_and\_not\_exists\_BKMData\_equalAmplitudeAntiProfileInitialVelocity\_affineAddScalarSchwartzPressure\_of\_exists\_affineCoeff\_ne\_zero\_on\_or\_active\_schwartzPressureGradient\_ne\_zero}.
It keeps the repaired finite-energy BKM clause available for
\(\texttt{twoModeSchwartzInitialVelocity}\ (a\,0)\ (a\,0)\ f\ (-f)\), but
rules out the matching structured BKM-data package over the same initial datum
for the bad affine-plus-localized pressure family.
The adjacent constructor-input repair is blocked too:
\texttt{ExplicitFiniteEnergyBKMContinuationClause\_equalAmplitudeAntiProfileInitialVelocity\_and\_not\_momentumEquation\_equalAmplitudeAntiProfileSchwartzVelocity\_affineAddScalarSchwartzPressure\_of\_exists\_affineCoeff\_ne\_zero\_or\_active\_schwartzPressureGradient\_ne\_zero}
keeps the repaired finite-energy BKM clause available for the cancelled
anti-profile initial datum, but proves that the full pointwise momentum
equation required by the finite-mode BKM clause constructor is impossible for
the same bad pressure family.
The clause normalization itself is now exposed in
\texttt{NavierStokesFiniteModeBKMClause.lean}:
\texttt{ExplicitBKMContinuationClause\_equalAmplitudeAntiProfileInitialVelocity\_iff\_zero}
and
\texttt{ExplicitFiniteEnergyBKMContinuationClause\_equalAmplitudeAntiProfileInitialVelocity\_iff\_zero}
state that the unrepaired and repaired BKM clauses for an equal-amplitude
anti-profile initial slice are exactly the zero-datum clauses.
Lean now closes the same cancelled anti-profile branch at the exact
whole-space output surface in
\texttt{NavierStokesFiniteModeAntiProfileGlobalOutputObstruction.lean}.  The
reusable theorem
\texttt{ExplicitConcreteNavierStokesGlobalOutputWithVelocityPressure\_zero\_iff\_spatialPressureGradient\_zero}
shows that a fixed zero-velocity/zero-datum global output exists exactly when
the proposed pressure has zero spatial gradient everywhere.  After applying
the anti-profile cancellation lemmas, this becomes
\texttt{ExplicitConcreteNavierStokesGlobalOutputWithVelocityPressure\_equalAmplitudeAntiProfileInitialVelocity\_velocity\_pressure\_iff\_spatialPressureGradient\_zero}:
an arbitrary smooth pressure can realize the cancelled equal-amplitude
anti-profile branch globally only in the spatially constant case.  Its direct
negative form
\texttt{not\_ExplicitConcreteNavierStokesGlobalOutputWithVelocityPressure\_equalAmplitudeAntiProfileInitialVelocity\_velocity\_pressure\_of\_exists\_spatialPressureGradient\_ne\_zero}
therefore lifts the earlier slabwise witness/BKM no-go to the exact whole-space
surface without committing to any pressure ansatz.  Lean also records the
matching affine-plus-localized specialization
\texttt{ExplicitConcreteNavierStokesGlobalOutputWithVelocityPressure\_equalAmplitudeAntiProfileInitialVelocity\_velocity\_affineAddScalarSchwartzPressure\_iff\_affineCoeff\_zero\_and\_schwartzPressureGradient\_zero\_if\_nonzeroAmplitude},
so even at the exact global-output layer the structured pressure family adds no
extra repair freedom beyond the already-collapsed zero-flow conditions.

That same finite-mode vorticity control now propagates through the corrected
energy/BKM bridge as well.  In
\texttt{NavierStokesEnergyBKMBridge.lean}, the theorems
\texttt{ExplicitFiniteEnergyBKMContinuationClause\_apply\_of\_add\_scalar\_smul\_schwartz\_of\_abs\_le\_of\_spatialDivergence\_zero\_affine\_add\_scalarSchwartzPressure},
\texttt{ExplicitFiniteEnergyBKMContinuationTarget\_apply\_of\_add\_scalar\_smul\_schwartz\_of\_abs\_le\_of\_spatialDivergence\_zero\_affine\_add\_scalarSchwartzPressure},
and their restricted-horizon companion remove the previously separate uniform
vorticity assumption from the bounded two-profile Schwartz energy route.  So
once the repaired finite-energy BKM clause (or target) itself is granted, the
existing energy identity, bounded-amplitude hypotheses, and Schwartz profiles
already suffice to drive this finite-mode lane all the way to the explicit
global-output conclusion.

There is now also a clean converse packaging for the exact bounded two-mode
closure route itself.  In \texttt{NavierStokesFiniteModeBKMClause.lean}, the
theorems
\texttt{ExplicitBKMContinuationClause\_twoModeSchwartzInitialVelocity\_of\_momentumEquation}
and
\texttt{ExplicitFiniteEnergyBKMContinuationClause\_twoModeSchwartzInitialVelocity\_of\_momentumEquation}
show that whenever the bounded two-mode Schwartz ansatz satisfies the full
pointwise momentum equation, the whole BKM clause family already holds on every
horizon for that datum.  So on this exact finite-mode lane, the BKM surface is
no longer merely a premise shell: it is fully inhabited because the earlier
finite-mode bounded-energy theorem has already produced whole-space explicit
global output.

At the same time, Lean now records the same theorem-shape bug on the unrepaired
BKM shell that had already been exposed on the older uniform-vorticity layer.
In \texttt{NavierStokesBKMContinuationTarget.lean}, theorem
\texttt{not\_ExplicitBKMContinuationTarget} shows that the raw BKM target is
false as stated because it still quantifies over negative horizons: on
\(T<0\), nonzero linear shear carries a formal finite-time witness and a
vacuous BKM envelope on the empty slab, while the demanded whole-space global
output is already impossible.  So only the repaired finite-energy or otherwise
nonnegative-horizon BKM formulations remain analytically meaningful.

The same exact closure route now also inhabits the older uniform-vorticity
continuation surface.  In
\texttt{NavierStokesFiniteModeUniformClause.lean}, the theorems
\texttt{ExplicitUniformVorticityContinuationClause\_twoModeSchwartzInitialVelocity\_of\_momentumEquation}
and
\texttt{ExplicitFiniteEnergyUniformVorticityContinuationClause\_twoModeSchwartzInitialVelocity\_of\_momentumEquation}
show that the bounded two-mode momentum closure already forces the unrepaired
and repaired uniform-vorticity clause families on every horizon as well.  So
for this exact finite-mode lane, both concrete continuation shells now sit
strictly below the already-proved whole-space explicit output theorem.

Lean now also isolates the missing finite-frame constant that turns a
componentwise curl estimate into the abstract \texttt{gamma}-bound used by the
Cole--Hopf identity interface.  In \texttt{IdentityEnergy.lean}, theorems
\texttt{gamma\_le\_card\_mul\_sq\_of\_abs\_le},
\texttt{abs\_frameEval\_le\_of\_abs\_le}, and
\texttt{abs\_vorticity\_le\_of\_componentwise\_abs\_le} show that if every
frame component is bounded in absolute value by \(M\), then the whole frame
energy is bounded by \((\#\iota) M^2\), hence the reconstructed vorticity is
bounded by the same explicit finite-cardinality factor.  Then
\texttt{ColeHopfInterface.lean} packages this as
\texttt{ColeHopfIdentityData.abs\_vorticity\_le\_of\_componentwise\_abs\_le},
so the abstract \texttt{curl\_energy} hypothesis can now be bypassed whenever
a concrete frame family comes with direct componentwise bounds.  This is real
progress, but it is deliberately limited: it does not yet produce those bounds
for the manuscript's actual frame families.

That finite-frame constant now reaches a concrete manuscript-facing package as
well.  In \texttt{GeometricColeHopfHeatApproximation.lean}, theorem
\texttt{geometricColeHopfHeatApproximation\_of\_componentwise\_abs\_le}
rebuilds the heat-decayed shared approximation shell directly from a
componentwise frame bound \(|\mathrm{curlFrame}_i(x)| \le M\), using the
derived explicit energy constant \((\#\iota) M^2\).  Then
\texttt{WindowedColeHopfHeatPackage.lean} specializes this to the genuine
smooth window cutoff via
\texttt{windowedColeHopfHeatSharedPackage\_of\_componentwise\_abs\_le} and
\texttt{windowedColeHopfHeat\_abs\_vorticity\_le\_uniform\_of\_componentwise\_abs\_le}.
So the current concrete windowed heat package no longer needs an abstract
\texttt{hcurl : gamma(\ldots) <= curlBound} input: a direct componentwise curl
bound is now sufficient to instantiate the uniform-vorticity tendril.  What
still remains open is deriving such a componentwise bound for the actual
manuscript frame family.

That same direct-bound route now reaches the first finite nonlocal shell too.
In \texttt{WindowedColeHopfHeatShiftKernelFrontier.lean}, the componentwise
counterparts
\texttt{windowedColeHopfHeatShiftKernelInitialSlice\_of\_componentwise\_abs\_le},
\texttt{windowedColeHopfHeatShiftKernelCandidate\_of\_componentwise\_abs\_le},
\texttt{toWindowedColeHopfHeatShiftKernelBridge\_of\_componentwise\_abs\_le},
and
\texttt{windowedColeHopfHeat\_realizes\_shiftKernel\_pressure\_seeded\_clause\_of\_componentwise\_abs\_le}
show that the concrete finite shift-kernel pressure-seeded clause can already
be instantiated from a uniform componentwise frame bound, not only from an
abstract \texttt{gamma}-energy hypothesis.  So the first genuine finite
translation-kernel route now sits on the same lowered assumption surface as the
underlying windowed package.

That same componentwise route now reaches the current clause-lift surface as
well.  In \texttt{GeometricColeHopfHeatFeffermanInterface.lean}, the
componentwise counterparts
\texttt{geometricColeHopfHeatSharedPackage\_of\_componentwise\_abs\_le},
\texttt{geometricColeHopfHeatUniformVorticityTendril\_of\_componentwise\_abs\_le},
\texttt{geometricColeHopfHeat\_abs\_vorticity\_le\_uniform\_of\_componentwise\_abs\_le},
and
\texttt{GeometricColeHopfHeatLiftData\_of\_componentwise\_abs\_le}
show that the present heat-decayed Fefferman interface can also be fed by a
direct componentwise frame bound, not only by an abstract \texttt{gamma}-bound.
So the remaining obstruction is no longer at the level of the vorticity
package or the current lift interface; it is the genuinely analytic task of
deriving a real frame bound and then supplying the missing velocity-pressure
lift data.

That same direct-bound route now reaches the current nonlinear saturated
pressure-seeded clause as well.  In
\texttt{WindowedColeHopfHeatSaturatedFrontier.lean}, the componentwise
counterparts
\texttt{windowedColeHopfHeatSaturatedInitialSlice\_of\_componentwise\_abs\_le},
\texttt{windowedColeHopfHeatSaturatedCandidate\_of\_componentwise\_abs\_le},
\texttt{toWindowedColeHopfHeatSaturatedBridge\_of\_componentwise\_abs\_le},
and
\texttt{windowedColeHopfHeat\_realizes\_saturated\_pressure\_seeded\_clause\_of\_componentwise\_abs\_le}
show that the present nonlinear saturated windowed route can already be fed by
a direct bound \(|\mathrm{curlFrame}_i(x)| \le M\), not only by an abstract
\texttt{gamma}-energy hypothesis.  This is a real lowering of the assumption
surface for the current concrete nonlinear clause.  What still remains is to
derive such a bound for an actual manuscript frame family, and then thread the
same replacement through the richer shift/sample frontiers and obstruction
files.

That same direct-bound route now reaches the generic nonlinear
shift-invariant saturated corridor as well.  In
\texttt{WindowedColeHopfHeatSaturatedShiftKernelInvariantFrontier.lean}, the
componentwise counterparts
\texttt{windowedColeHopfHeatSaturatedShiftKernelInitialSlice\_eq\_saturatedInitialSlice\_of\_shiftInvariant\_constantMagnitude\_of\_componentwise\_abs\_le},
\texttt{windowedColeHopfHeatSaturatedShiftKernelCandidate\_eq\_saturatedCandidate\_of\_stationary\_shiftInvariant\_constantMagnitude\_of\_componentwise\_abs\_le},
\texttt{windowedColeHopfHeat\_realizes\_saturated\_pressure\_seeded\_clause\_via\_shiftKernel\_of\_shiftInvariant\_constantMagnitude\_of\_componentwise\_abs\_le},
and
\texttt{exists\_windowedColeHopfHeatSaturatedShiftKernelBridge\_of\_stationary\_shiftInvariant\_constantMagnitude\_of\_componentwise\_abs\_le}
show that the broader saturated shift-kernel invariant bridge now also
instantiates from a direct componentwise frame bound, not only from an
abstract \texttt{gamma}-energy hypothesis.  So the remaining shift-kernel-side
burden is no longer to lower the invariant frontier itself; it is to push the
same replacement through the obstruction files and the richer masked/anchored
specializations that probe this corridor.

That next obstruction step is now partially complete too.  In
\texttt{WindowedColeHopfHeatSaturatedShiftKernelInvariantObstruction.lean}, the
componentwise counterparts
\texttt{windowedColeHopfHeatShiftKernel\_abs\_eq\_abs\_of\_topDownBridge\_eq\_saturatedCandidate\_of\_stationary\_shiftInvariant\_of\_componentwise\_abs\_le},
\texttt{windowedColeHopfHeatShiftKernel\_totalCoeffSum\_zero\_of\_topDownBridge\_eq\_saturatedCandidate\_of\_abs\_ne\_abs\_of\_stationary\_shiftInvariant\_of\_componentwise\_abs\_le},
and
\texttt{not\_exists\_windowedColeHopfHeatShiftKernelTopDownBridge\_eq\_saturatedCandidate\_of\_abs\_ne\_abs\_of\_stationary\_shiftInvariant\_of\_componentwise\_abs\_le}
show that the bridge-level obstruction itself also descends to the concrete
componentwise frame-bound surface.  So the present invariant shift-kernel lane
is now lowered on both sides: its positive bridge package and its decisive
nonexistence witness are both stated without any separate abstract
\texttt{gamma}-energy hypothesis.  The remaining burden is to propagate that
same lowering through the richer masked/anchored descendants and the more
specialized obstruction shells built on top of them.

The bridge-level coefficient-size obstruction has now been lowered as well.  In
the same file, the componentwise counterparts
\texttt{windowedColeHopfHeatShiftKernel\_abs\_totalCoeffSum\_le\_abs\_saturatedCoeff\_of\_topDownBridge\_eq\_saturatedCandidate\_of\_stationary\_shiftInvariant\_of\_componentwise\_abs\_le},
\texttt{windowedColeHopfHeatShiftKernel\_abs\_totalCoeffSum\_lt\_abs\_saturatedCoeff\_of\_topDownBridge\_eq\_saturatedCandidate\_of\_totalCoeffSum\_ne\_zero\_of\_stationary\_shiftInvariant\_of\_componentwise\_abs\_le},
\texttt{not\_exists\_windowedColeHopfHeatShiftKernelTopDownBridge\_eq\_saturatedCandidate\_of\_abs\_saturatedCoeff\_lt\_abs\_totalCoeffSum\_of\_stationary\_shiftInvariant\_of\_componentwise\_abs\_le},
and
\texttt{not\_exists\_windowedColeHopfHeatShiftKernelTopDownBridge\_eq\_saturatedCandidate\_of\_abs\_saturatedCoeff\_le\_abs\_totalCoeffSum\_of\_totalCoeffSum\_ne\_zero\_of\_stationary\_shiftInvariant\_of\_componentwise\_abs\_le}
show that the aggregate coefficient-mass contradiction now also lives on the
same explicit finite-frame bound surface.  So the generic invariant
shift-kernel corridor now has both of its current bridge-level obstruction
families on concrete componentwise hypotheses: the varying-magnitude witness
route and the coefficient-size route.

The first anchored descendant now reaches that same concrete surface on its
constructive side.  In
\texttt{WindowedColeHopfHeatShiftKernelFrontier.lean}, Lean now exposes the
anchored wrappers
\texttt{windowedColeHopfHeatAnchoredShiftInitialSlice\_of\_componentwise\_abs\_le},
\texttt{windowedColeHopfHeatAnchoredShiftCandidate\_of\_componentwise\_abs\_le},
\texttt{windowedColeHopfHeatAnchoredTranslationShiftKernelInitialSlice\_eq\_anchoredShiftInitialSlice\_of\_componentwise\_abs\_le},
and
\texttt{windowedColeHopfHeatAnchoredTranslationShiftKernelCandidate\_eq\_anchoredShiftCandidate\_of\_componentwise\_abs\_le}.
Then \texttt{WindowedColeHopfHeatSaturatedAnchoredShiftInvariantFrontier.lean}
adds the componentwise counterparts
\texttt{windowedColeHopfHeatAnchoredShiftCandidate\_eq\_saturatedCandidate\_of\_stationary\_shiftInvariant\_constantMagnitude\_of\_componentwise\_abs\_le},
\texttt{windowedColeHopfHeat\_realizes\_saturated\_pressure\_seeded\_clause\_via\_anchoredShift\_of\_shiftInvariant\_constantMagnitude\_of\_componentwise\_abs\_le},
and
\texttt{exists\_windowedColeHopfHeatSaturatedAnchoredShiftKernelBridge\_of\_stationary\_shiftInvariant\_constantMagnitude\_of\_componentwise\_abs\_le}.
So the anchored shift-invariant saturated corridor no longer needs a separate
abstract \texttt{gamma}-energy bound on its positive bridge-construction side.
What still remains there is the obstruction side, which still sits above this
surface.

That obstruction side is now lowered as well for the first anchored shell.  In
\texttt{WindowedColeHopfHeatSaturatedAnchoredShiftInvariantObstruction.lean},
Lean now proves the componentwise candidate-side theorems
\texttt{windowedColeHopfHeatAnchoredShift\_abs\_eq\_abs\_of\_candidate\_eq\_saturatedCandidate\_of\_stationary\_shiftInvariant\_of\_componentwise\_abs\_le},
\texttt{windowedColeHopfHeatAnchoredShift\_abs\_coeffSum\_lt\_abs\_saturatedCoeff\_of\_candidate\_eq\_saturatedCandidate\_of\_stationary\_shiftInvariant\_of\_componentwise\_abs\_le},
\texttt{not\_windowedColeHopfHeatAnchoredShiftCandidate\_eq\_saturatedCandidate\_of\_abs\_saturatedCoeff\_le\_abs\_coeffSum\_of\_stationary\_shiftInvariant\_of\_componentwise\_abs\_le},
and
\texttt{not\_windowedColeHopfHeatAnchoredShiftCandidate\_eq\_saturatedCandidate\_of\_abs\_ne\_abs\_of\_stationary\_shiftInvariant\_of\_componentwise\_abs\_le},
together with the matching bridge-level counterparts
\texttt{windowedColeHopfHeatAnchoredShift\_abs\_eq\_abs\_of\_topDownBridge\_eq\_saturatedCandidate\_of\_stationary\_shiftInvariant\_of\_componentwise\_abs\_le},
\texttt{windowedColeHopfHeatAnchoredShift\_abs\_coeffSum\_lt\_abs\_saturatedCoeff\_of\_topDownBridge\_eq\_saturatedCandidate\_of\_stationary\_shiftInvariant\_of\_componentwise\_abs\_le},
\texttt{not\_exists\_windowedColeHopfHeatAnchoredShiftTopDownBridge\_eq\_saturatedCandidate\_of\_abs\_saturatedCoeff\_le\_abs\_coeffSum\_of\_stationary\_shiftInvariant\_of\_componentwise\_abs\_le},
and
\texttt{not\_exists\_windowedColeHopfHeatAnchoredShiftTopDownBridge\_eq\_saturatedCandidate\_of\_abs\_ne\_abs\_of\_stationary\_shiftInvariant\_of\_componentwise\_abs\_le}.
So the first anchored shift-invariant obstruction family now also lives
directly on the explicit componentwise frame-bound surface.  The remaining
anchored burden is no longer this basic obstruction shell; it is the richer
masked and downstream descendant files.

The specialization layer tying that older anchored exactness theorem back to
the newer full invariant finite-kernel theorem is now lowered to the same
surface as well.  In
\texttt{WindowedColeHopfHeatSaturatedShiftKernelInvariantAnchoredSpecialization.lean},
the new theorems
\texttt{windowedColeHopfHeatAnchoredShift\_abs\_eq\_abs\_of\_candidate\_eq\_saturatedCandidate\_of\_stationary\_shiftInvariant\_via\_shiftKernelInvariant\_of\_componentwise\_abs\_le}
and
\texttt{not\_windowedColeHopfHeatAnchoredShiftCandidate\_eq\_saturatedCandidate\_of\_abs\_ne\_abs\_of\_stationary\_shiftInvariant\_via\_shiftKernelInvariant\_of\_componentwise\_abs\_le}
show that the anchored candidate-side exactness and impossibility statements
are now not only proved directly in the anchored file, but already recovered as
specializations of the broader invariant shift-kernel theorem on that same
concrete componentwise frame-bound surface.

\section{What Still Needs Heavy Analysis}

The new Lean kernel also clarifies what it does \emph{not} solve.

\begin{enumerate}[leftmargin=1.5em]
  \item It does not prove the local propagation estimate for
  \(\Gamma(P_t \phi)(\mathrm{id})\).
  \item It does not prove the manuscript-shaped truncation theorem that would
  feed convergence of the radius and statistic observables into the new
  approximation shell, or the eventual positivity / coefficient-energy bounds
  into the new finite-mode limit shell.
  \item It does not yet lift the current concrete windowed candidate and its
  descendants to real velocity-generation or pressure-generation theorems
  beyond the present seeded local predicates.
  \item It does not turn time-independent compatibility seeds into local
  solutions; Lean now records that this requires the stationary momentum
  balance, or else a genuinely evolved velocity-pressure construction.
  \item It does not justify the manuscript's push-down theorem at a fully
  analytic infinite-dimensional level.
  \item It does not close BKM in general.  The new Lean theorem closes the BKM
  premise only for the bounded two-mode Schwartz ansatz; it does not yet supply
  the corresponding analytic continuation input for the manuscript's broader
  infinite-dimensional routes.
\end{enumerate}

That honesty is a feature, not a weakness.  The point of grassroots work is to
shrink the unknown region and expose the remaining burden sharply.

\section{Next Grassroots Targets}

The next constructive steps are fairly clear.

\begin{enumerate}[leftmargin=1.5em]
  \item Push the concrete windowed route through its first real compatibility
  theorem.  The present fixed candidate, scaled descendant, and affine
  seed-blend descendant already have their local certificates; the next theorem
  should decide whether any of them can generate a nontrivial velocity field in
  a way that matches the current vorticity tendril.
  \item Strengthen the new local kernel-founded shell from ``kernel semigroup +
  identity interface'' to ``kernel semigroup + positivity/energy propagation''.
  That remains the next honest place where a real analytic theorem can plug in.
  \item Replace the current selected-coefficient heat surrogate by a more
  faithful finite-mode identity differential law, so the concrete shared route
  is closer to the manuscript's actual SG/identity story.
  \item Derive a genuine componentwise curl bound for one of the manuscript's
  actual explicit frame families, so the new
  \texttt{windowedColeHopfHeatSharedPackage\_of\_componentwise\_abs\_le}
  constructor can be fed by internal route data rather than an external
  hypothesis.
\end{enumerate}

The reason for this order is straightforward: land a small exact kernel before
spending effort on a speculative infinite-dimensional closure.

\section{Conclusion}

The SG-Cole-Hopf route is not formally closed.  But it is already structured
enough that its algebraic core can be extracted and verified independently.

That is what the current grassroots pass accomplishes.  The new Lean module does
not prove global regularity, but it does prove that the manuscript's
identity-energy route has a small reusable spine:
\[
  \Gamma(P_t \phi)(\mathrm{id})
  \Longrightarrow
  \Gamma(W(t))(\mathrm{id})
  \Longrightarrow
  |\omega(t,x)|
\]
once the hard analytic hypotheses are fed in.

That is exactly the kind of background theory we want from the ground up.

\end{document}
